% arara: pdflatex
% arara: biber
% arara: pdflatex
% arara: pdflatex

\documentclass[smallextended]{svjour3}       % onecolumn (second format)
\smartqed  % flush right qed marks, e.g. at end of proof
\usepackage{graphicx}

% ---------- Math & symbols ----------
\usepackage{amsmath,amssymb,bm,mathtools}
\usepackage{stmaryrd}   % \llbracket,\rrbracket

% ---------- Figures & tables ----------
\usepackage{graphicx}
\usepackage{subcaption}
\usepackage{tabularx}
\usepackage{booktabs}
\usepackage{array,longtable}
\setlength{\LTcapwidth}{\textwidth} % longtable captions full width

% ---------- References & hyperlinks ----------
\usepackage{xcolor}
\usepackage{hyperref}
\hypersetup{
	colorlinks=true,
	linkcolor=blue,
	citecolor=violet,
	urlcolor=teal
}

\usepackage[backend=biber,style=numeric,sorting=none]{biblatex}
\addbibresource{Selivanov.bib}

\journalname{International Journal of Fracture}

% ===================== Notation layer  =====================

% --- Vectors & 2nd-order tensors: bold italic ---
\newcommand{\vb}[1]{\bm{#1}}      % generic vector-like (also 2nd-order tensors)
%\newcommand{\tens}[1]{\bm{#1}}    % alias for clarity in code

% --- Keep access to arrowed vectors; make \vec bold-italic by default ---
\let\arrowvec\vec                  % save the original arrowed \vec
\renewcommand{\vec}[1]{\bm{#1}}    % now \vec{x} -> bold-italic x

% --- Upright bold matrices / discrete operators ---
\newcommand{\mat}[1]{\mathbf{#1}}  % e.g., \mat{K}, \mat{M}

% --- 4th-order tensors (blackboard) ---
\newcommand{\four}[1]{\mathbb{#1}} % e.g., \four{D}

% --- Unit vectors & local frame ---
\newcommand{\ex}{\vb{e}_x}
\newcommand{\ey}{\vb{e}_y}
\newcommand{\tvec}{\vb{t}}
\newcommand{\nvec}{\vb{n}}
\newcommand{\QQ}{\mat{Q}}          % Leg-2 local frame matrix [t^T; n^T]
% (No conflict with \mathbb{Q}; that remains available as is.)

% --- Common fields ---
\newcommand{\uu}{\vb{u}}
\newcommand{\eps}{\tens{\varepsilon}}
\newcommand{\sig}{\tens{\sigma}}

% --- Jumps & averages (starred version auto-sizes) ---
\DeclarePairedDelimiter{\jump}{\llbracket}{\rrbracket}
\DeclarePairedDelimiter{\avg}{\langle}{\rangle}
% Usage: \jump{u}, \jump*{u};  \avg{\cdot}, \avg*{\cdot}

% --- Other handy paired delimiters ---
\DeclarePairedDelimiter{\abs}{\lvert}{\rvert}
\DeclarePairedDelimiter{\norm}{\lVert}{\rVert}

% --- Operators (proper spacing in math) ---
\DeclareMathOperator{\sgn}{sgn}
\DeclareMathOperator{\tr}{tr}
\DeclareMathOperator{\sym}{sym}
\DeclareMathOperator{\skw}{skw}
\newcommand{\diver}{\operatorname{div}}
\newcommand{\grad}{\nabla}

% --- Differential and helpers ---
\newcommand{\dd}{\mathop{}\!\mathrm{d}}  % correct spacing: \int f\,\dd x
\newcommand{\nrm}[1]{\overline{#1}}      % normalized quantities (bar)

% --- Backward-compatibility (new-or-renew without @) ---
% \providecommand defines if undefined; \renewcommand then overwrites if needed.
\providecommand{\bt}{\tvec}
\renewcommand{\bt}{\tvec}

\providecommand{\bn}{\nvec}
\renewcommand{\bn}{\nvec}

\providecommand{\bmU}{\uu}
\renewcommand{\bmU}{\uu}

\providecommand{\bd}{\vb d}
\renewcommand{\bd}{\vb d}

\providecommand{\bTcz}{\vb{T}_{\mathrm{cz}}}
\renewcommand{\bTcz}{\vb{T}_{\mathrm{cz}}}

\providecommand{\be}{\vb e}
\renewcommand{\be}{\vb e}

\providecommand{\DD}{\four D}
\renewcommand{\DD}{\four D}

\providecommand{\mm}{\QQ}
\renewcommand{\mm}{\QQ}
% =================== End Notation layer ====================

	\include{changes.tex}

\begin{document}



\title{Potential-Based Cohesive Modeling of Mixed-Mode Crack Kinking.}

\author{Viacheslav Bogdanov          \and
	Danylo Selivanov %etc.
}

%\authorrunning{Short form of author list} % if too long for running head

\institute{Viacheslav Bogdanov \at
	S.P. Timoshenko Institute of Mechanics, NAS of Ukraine, 3 Nesterov Street, Kyiv, 03057, Ukraine \\
	\email{bogdanov@nas.gov.ua}   \\
	\and
	Danylo Selivanov \at
	S.P. Timoshenko Institute of Mechanics, NAS of Ukraine, 3 Nesterov Street, Kyiv, 03057, Ukraine
	\\
	\email{fract@inmech.kyiv.ua}
}

\date{Received: date / Accepted: date}
% The correct dates will be entered by the editor


\maketitle

\begin{abstract}
	We present a compact finite-element framework for crack kinking in plane strain that augments the physical crack with a short, zero-thickness cohesive continuation. A thin auxiliary channel is used to generate a refined T3 mesh and is collapsed after T6 conversion, so the analysis is performed on the physical interface with coincident node pairs. The interface law is potential-based, differentiable, and bounded in both normal and tangential directions; its cohesive potential yields an energetically consistent formulation with a symmetric Hessian tangent, enabling robust Newton convergence. Local openings are evaluated in a frame attached to the trial leg and rotated to the global frame for assembly. We detail the construction of interface degrees of freedom, the jump operator and rotation, and a mode-mix-aware procedure to select the kink direction.
	Numerical studies show stable determination of the kink angle, rapid convergence of the critical load, and reproducible energy measures across independently generated meshes. With moderate far-field caps and interface resolution, the critical load is mesh-independent to near machine precision, and bulk tractions on the faces match cohesive tractions reconstructed from jumps along almost the entire leg. We also compare a global load-based selector with a local energetic mouth criterion, highlighting their complementary behavior. The framework is designed as a practical, mesh-robust tool for simulating initiation and early growth of kinking cracks under mixed-mode conditions.
	\keywords{Mixed-mode fracture \and Cohesive zone model \and Traction--separation law \and Potential-based cohesive interface}
	% \PACS{PACS code1 \and PACS code2 \and more}
	%\subclass{MSC code1 \and MSC code2 \and more}
\end{abstract}

\section{Introduction}

Accurate prediction of how and where cracks deflect is essential in many
engineering systems, from layered composites and adhesive joints to
microelectronic packages and structural components subjected to non-uniform
loading. Even small deviations in crack trajectory can alter load-carrying
capacity, trigger delamination, or precipitate catastrophic failure.
Classical Linear Elastic Fracture Mechanics (LEFM) provides canonical
predictions of the incipient kink direction, notably via the maximum
circumferential-stress criterion and related energy-based arguments
\cite{ErdoganSih1963}. Path conditions such as enforcing
$K_{\rm II}=0$ at an advancing crack tip \cite{CotterellRice1980}
provide equivalent selection principles. For cracks interacting with
material interfaces, early analyses established when a crack continues
along the interface or kinks into an adjacent medium, highlighting the
dominant role of the mixed-mode elastic field in steering the initial
deflection \cite{HeHutchinson1989}.

Cohesive-zone formulations complement LEFM by embedding fracture processes
directly in traction--separation relations (TSLs), thereby linking intact
material response with complete decohesion without explicitly tracking the
crack. Foundational models idealised the process zone using simple cohesive
laws \cite{Barenblatt1962,Dugdale1960}, and subsequent finite-element
implementations showed that such laws can robustly capture initiation,
deflection, branching, and arrest in brittle and quasi-brittle materials
\cite{CamachoOrtiz1996}. Cohesive laws have also been central to
interpreting and predicting resistance curves
\cite{TvergaardHutchinson1992} and to modelling dynamic crack growth with
intrinsic interface behaviour \cite{XuNeedleman1994}.

\paragraph{Mixed-mode fracture and cohesive laws.}
A comprehensive review by Park and Paulino \cite{Park2011Review}
distinguishes two broad families of mixed-mode models:
\emph{energy-based} criteria and \emph{traction-based} cohesive laws.
Energy-based schemes---such as the widely used Benzeggagh--Kenane (BK)
criterion \cite{BenzeggaghKenane1996}---prescribe a mode-dependent
fracture toughness but do not supply interface tractions and cannot be
derived from a single potential. Traction-based cohesive laws, initiated
by Xu and Needleman \cite{XuNeedleman1994}, define an explicit mixed-mode
potential $\Psi(d_n,d_t)$ whose derivatives give the physical tractions.
Subsequent extensions introduced polynomial, power-law, exponential, and
asymmetric couplings, allowing the mode transition and shear sensitivity
to be tuned to the material system. Such potential-based cohesive laws
ensure thermodynamic consistency and yield symmetric algorithmic tangents,
which is advantageous for implicit finite-element formulations and
particularly relevant in mixed-mode initiation and kinking problems.

The importance of accurately describing the near-tip fracture-process zone
(FPZ) has motivated several recent modelling developments. Najmeddine,
Gupta and Moini \cite{NajmeddineGuptaMoini2025} coupled a phase-field
description of bulk fracture with a PPR cohesive interface to capture
deflection, penetration and toughening mechanisms in multi-material
systems, while Shabrou and Daneshjoo \cite{ShabrouDaneshjoo2025} proposed
mixed-mode criteria for orthotropic materials that incorporate FPZ effects
and T-stress, improving directional predictions over classical LEFM
criteria. Analytical progress has likewise sharpened the understanding of
mixed-mode fields: Liu and Wei derived a closed-form stress field for
finitely kinked cracks \cite{Liu2021KinkedStressFields} and later a
semi-analytical stress-intensity-factor framework for branched cracks
\cite{Liu2023BranchedSIF}, Si and Wei extended these ideas to kinked edge
cracks in a half-space \cite{Si2024KinkedEdge}, and Shi and Xu obtained
explicit stress-intensity factors for infinitesimal kinks in orthotropic
solids \cite{Shi2025OrthotropicKinkSIF}. Parallel cohesive-zone studies
have clarified how mode mixity, cohesive strength and interface toughness
compete to determine crack deflection versus penetration
\cite{ParmigianiThouless2006CZMDeflection,Alam2016IncidentAngleCZM,
	Tankasala2017KinkingOrthotropic,ChenDeBorst2022PFRCZM}.

\paragraph{Motivating example: why cohesive models remain relevant even in a near-LEFM regime.}
Even in configurations where LEFM is nominally applicable, cohesive-zone
effects can produce measurable changes in the predicted critical load. For
a mode-I crack of half-length $a$ in an infinite plate, LEFM predicts
onset of fracture when $K_{\mathrm{I}}=\sigma\sqrt{\pi a}$ reaches the
toughness $K_{\mathrm{Ic}}$, giving
$\sigma_{\mathrm{cr}}^{\mathrm{LEFM}} = K_{\mathrm{Ic}}/\sqrt{\pi a}$.
A cohesive-zone representation such as the Dugdale model replaces the
singular elastic field by a finite process zone of length
$\ell_{\mathrm{p}}$ and consequently yields a slightly smaller critical
load. Even when the process-zone size is only about $5\%$ of the crack
half-length---a regime often treated in practice as close to the LEFM
limit---the cohesive prediction remains $2$--$3\%$ lower than the LEFM
estimate. This modest but systematic deviation shows that a finite process
zone influences not only the near-tip fields but also global quantities
such as the critical load.

This example motivates the viewpoint adopted in the present work. The
question is not whether LEFM should be replaced in the small-process-zone
regime, but how a regularized cohesive formulation can be used to assess,
in a controlled way, which parts of the kinking response remain essentially
elastic and which retain a measurable imprint of the traction--separation
law and of the finite FPZ. In particular, the present study is conducted
deliberately in a near-LEFM regime, so that agreement and discrepancy with
classical predictors can both be interpreted meaningfully.

\paragraph{Aim and modelling philosophy.}
Our aim is to isolate, in the cleanest possible setting, how mode mixity,
cohesive strength, and crack orientation conspire to determine the kink
direction under mixed-mode loading, and how this choice manifests in the
critical load. To this end, we employ a deliberately minimal plane-strain
finite-element formulation in which the physical crack is extended by a
short, zero-thickness cohesive continuation. A slender auxiliary channel
is introduced solely to control mesh resolution and is collapsed prior to
loading, ensuring that all tractions and boundary conditions apply
directly to the true crack faces. For each trial orientation, we compute a
critical state dictated solely by the prescribed traction--separation law
and examine two complementary selectors for the kink direction: a global
load-based criterion and a local cohesive-energetic criterion at the
mouth.

A second objective is to distinguish clearly between constitutive effects
and auxiliary numerical choices. To this end, the analysis examines how the
predicted direction depends on the finite trial cohesive length used to
regularize the near-tip field, and identifies the range of applicability
and the potential failure modes of the method. This point is important
because the present framework does not evaluate a tip criterion at a
mathematical crack tip; instead, it infers the preferred direction from a
short cohesive continuation of finite extent. The method should therefore
be interpreted as a finite-process-zone formulation intended for controlled
comparison with LEFM in a small but nonzero process-zone regime.

This modelling philosophy avoids remeshing, enrichment, or
geometry-specific heuristics and reveals, with high transparency, which
aspects of the cohesive law influence the kink direction and which enter
primarily through the evolution of the fracture-process zone and the
critical load. At the same time, it provides a natural setting for
examining numerical robustness, including the sensitivity to the auxiliary
cohesive-leg length and the conditioning of different directional
selectors.

\paragraph{Structure of the paper.}
Section~\ref{sec:problem} introduces the problem statement, geometry and
notation. Section~\ref{sec:governing} presents the governing equations of
the elastic bulk. Section~\ref{sec:tsl} develops the cohesive interface
formulation, mixed-mode potential, and consistent linearisation.
Section~\ref{sec:construction} describes interface operators and assembly.
Section~\ref{sec:mesh} outlines mesh construction, channel collapse, and
refinement. Section~\ref{sec:results} reports numerical results, including
field plots, convergence studies, and kink selection.
Section~\ref{sec:LEFM-comparison} compares cohesive predictions with LEFM
benchmarks. Section~\ref{sec:scope_limitations} examines the sensitivity
to the auxiliary cohesive-leg length and states the range of applicability
and limitations of the method. Section~\ref{sec:conclusions} summarises
the findings and future directions.

{\small
	\begin{longtable}{@{}>{\raggedright\arraybackslash}p{0.22\textwidth}%
			>{\raggedright\arraybackslash}p{0.64\textwidth}%
			>{\raggedleft\arraybackslash}p{0.08\textwidth}@{}}
		\caption{Notation used throughout. Subscripts $n,t$ denote normal and tangential components in the Leg~2 frame $\mat Q=[\bt^\top;\bn^\top]$.}
		\label{tab:notation}\\
		\toprule
		\textbf{Symbol} & \textbf{Meaning} & \textbf{Units}\\
		\midrule
		\endfirsthead
		\multicolumn{3}{@{}l@{}}{\small\textit{Table \thetable\ (continued)}}\\
		\toprule
		\textbf{Symbol} & \textbf{Meaning} & \textbf{Units}\\
		\midrule
		\endhead
		\midrule
		\multicolumn{3}{r}{\small\textit{Continued on next page}}\\
		\endfoot
		\bottomrule
		\endlastfoot

		\multicolumn{3}{@{}l@{}}{\textbf{Geometry and loading}}\\
		\midrule
		$A,\,B$ & Plate width and half-height ($x\in[0,A],\,y\in[-B,B]$) & m\\
		$V_0,\,V_1,\,V_2$ & Mouth, knee, and tip coordinates of the two-leg path & m\\
		$a$ & Physical leg length (Leg~1): $|V_1{-}V_0|$ & m\\
		$\delta$ & Prescribed cohesive length (Leg~2, trial leg) & m\\
		$\beta$ & Active cohesive sublength on Leg~2 & m\\
		$\theta_1,\,\theta_2$ & Absolute orientations of Leg~1/2 (from $+x$) & rad\\
		$\Delta\theta$ & Signed kink angle $\operatorname{wrap}_{(-\pi,\pi]}(\theta_2-\theta_1)$ & rad\\
		$\bt,\,\bn$ & Unit tangent (knee$\to$tip) and unit normal (lower$\to$upper) & --\\
		$\mat Q=[\bt^\top;\bn^\top]$ & Local $2{\times}2$ frame attached to Leg~2 & --\\
		$s,\,r$ & Local coordinates along $\bt$ and $\bn$ & m\\
		$\sigma$ & Magnitude of the applied tractions on the top and bottom plate faces (equal in magnitude, opposite in sign) & MPa\\[2pt]

		\multicolumn{3}{@{}l@{}}{\textbf{Bulk kinematics and material}}\\
		\midrule
		$\vec u$ & Displacement field & m\\
		$\jump{\vec u}$ & Displacement jump across Leg~2: $\vec u^{(+)}-\vec u^{(-)}$ & m\\
		$\boldsymbol\varepsilon(\vec u)$ & Small strain $\tfrac12(\nabla\vec u+\nabla\vec u^\top)$ & --\\
		$\boldsymbol\sigma$ & Cauchy stress & Pa\\
		$\mat D$ & Elasticity tensor (plane strain) & Pa\\
		$E,\ \nu,\ G$ & Young’s modulus, Poisson’s ratio, shear modulus ($G=E/[2(1+\nu)]$) & Pa, --, Pa\\[2pt]

		\multicolumn{3}{@{}l@{}}{\textbf{Mesh and interface}}\\
		\midrule
		$n_{\mathrm{coh}}$ & Equal segments on Leg~2 (interface resolution) & --\\
		$h_{\min}$ & Target near-interface size $=\delta/n_{\mathrm{coh}}$ & m\\
		$h_{\mathrm{grad}}$ & Geometric growth factor away from the band & --\\
		$h_{\max}$ & Far-field size cap & m\\
		$w$ & Channel half-width (mesh only; collapsed before solve) & m\\[2pt]

		\multicolumn{3}{@{}l@{}}{\textbf{Cohesive variables and TSL}}\\
		\midrule
		$\vec d=(d_t,d_n)$ & Local sliding/opening in $(\bt,\bn)$: $\vec d=\mat Q\,\jump{\vec u}$ & m\\
		$d_n^{\max},\,d_t^{\max}$ & Critical separations (normal, tangential) & m\\
		$\bar d_n,\,\bar d_t$ & Normalized separations: $d_n/d_n^{\max}$, $|d_t|/d_t^{\max}$ & --\\
		$\vec T_{\rm cz}=(\sigma_t,\sigma_n)$ & Local cohesive traction (action–reaction on faces) & Pa\\
		$\sigma_n^{\max},\,\sigma_t^{\max}$ & Cohesive strengths (normal, shear) & Pa\\
		$\bar\sigma_n,\,\bar\sigma_t$ & Normalized tractions $\sigma_\alpha/\sigma_\alpha^{\max}$ & --\\
		$\phi_n,\,\phi_t$ & Mode~I/II fracture energies (areas under TSL branches) & J\,m$^{-2}$\\
		$F,t,g$ & Normalized cohesive potential and its derivatives
		($t=F'$, $g=t'$), with saturation \chg{R2C3}{$F(1)=c$} & --\\
		\chg{R2C3}{$c$} & \chg{R2C3}{Dimensionless TSL shape parameter ($c=\int_0^1 t(\xi)\,\mathrm d\xi$)} & --\\
		\chg{R2C3}{$\chi_\alpha$} &
		\chg{R2C3}{Mode-wise energetic prefactor of the cohesive law,
		$\chi_\alpha=\sigma_\alpha^{\max} d_\alpha^{\max}
		=\phi_\alpha/c_\alpha$}
		&
		J\,m$^{-2}$\\
		$a_1,\,a_2$ & TSL shape switches (end of strengthening branch; onset of softening) & --\\
		
		$\kappa$ & Dimensionless compression penalty (e.g., $2/a_1$) & --\\
		$k_c$ & Contact slope $k_c=\kappa\,\sigma_n^{\max}/d_n^{\max}$ & Pa\,m$^{-1}$\\
		$r$ & Mixed-mode energetic coupling parameter ($0\le r\le r^{\max}$) & --\\[2pt]

		\multicolumn{3}{@{}l@{}}{\textbf{Discrete interface operators (assembly)}}\\
		\midrule
		$\mathcal T$ & Nodewise TSL map, $\vb s_e=\mathcal T(\vec d_e)$ & --\\
		$\mat R$ & Consistent line “mass’’: $\mat R=l_e\!\int_0^1\!N^\top N\,\mathrm d\xi$ & m\\
		$\mat A_e$ & Face assembly/rotation: $\vec f_e=\mat A_e\,\vb s_e$ (action–reaction) & --\\
		$\mat B_e$ &
		Element jump operator mapping nodal displacements to local displacement jumps,
		$\bd_e=\mat B_e\vb U_e$
		&
		--\\
		$\mat P_e$ & Node-order reversal $(i,j,k)\mapsto(k,j,i)$ & --\\
		$\mat K_{\rm loc}$ & Local $2\times2$ tangent & Pa\,m$^{-1}$\\[2pt]

		\multicolumn{3}{@{}l@{}}{\textbf{Outputs and selectors}}\\
		\midrule
		$\sigma_{\mathrm{cr}}$ & Critical load at the critical state & Pa\\
		$\Psi_{\mathrm{crit}},\ \hat\Psi_{\mathrm{crit}}$ & Mouth work on Leg~2; $\hat\Psi_{\mathrm{crit}}=\Psi_{\mathrm{crit}}/\phi_n$ & J\,m$^{-2}$, --\\
		$\theta_2^{\sigma},\,\theta_2^{\Psi}$ & Kink angles from load-based and mouth-energy criteria & deg\\
		$\sigma_\infty$ & Extrapolated critical load (Richardson/Aitken) & Pa\\
		$p,\ \epsilon$ & Observed order and finest-mesh relative error & --\\
	\end{longtable}
}


\section{Problem statement}
\label{sec:problem}

We consider a two-dimensional body in plane strain containing a predefined
two-leg crack path that serves as the basis for studying mixed-mode crack
kinking under quasi-static loading conditions. The problem setup, geometry,
local frames, and representative traction--separation law are summarized in
Fig.~\ref{fig:statement}.

The body, shown in Fig.~\ref{fig:stmt:a}, is a rectangular domain subjected
to uniform tensile tractions on its top and bottom edges, while the vertical
sides remain free of external loads. The internal discontinuity consists of
two consecutive legs forming a polyline. The first leg (Leg~1) represents the
existing physical crack of prescribed length and orientation, whereas the
second leg (Leg~2) is a short fictitious continuation introduced to realize a
cohesive interface on which kinking may occur. We denote by
$\Gamma_{\mathrm{cr}}$ the boundary of the physical crack (Leg~1), and by
$\Gamma_{\mathrm{cz}}$ the boundary of the cohesive extension (Leg~2). The
kink originates at the common junction point of the two legs, referred to
here as the \emph{knee}, and terminates at the trial cohesive tip. The
orientation of Leg~2 relative to Leg~1 defines the kink angle. The material
is assumed homogeneous and isotropic, and plane-strain conditions are imposed
throughout.

A closer view of the two-leg configuration is presented in
Fig.~\ref{fig:stmt:b}. The physical crack, Leg~1, is traction-free, whereas
Leg~2 carries cohesive tractions acting on its two opposing faces. These
tractions depend on the local relative displacements between the upper and
lower faces and are expressed in a local coordinate system rigidly attached
to Leg~2. This local frame consists of a tangential direction, aligned with
the cohesive segment, and a normal direction, pointing from the lower to the
upper face. The corresponding tractions therefore form a pair of equal and
opposite vectors acting on the two sides of the cohesive zone.

The distribution of openings and tractions along the cohesive continuation is
sketched in Fig.~\ref{fig:stmt:c}. Along Leg~2, the normal opening and the
tangential sliding evolve according to the adopted traction--separation law.
The active cohesive region is bounded by the physical crack tip and the trial
cohesive tip, whereas the remaining crack faces are traction-free. The local
coordinates $(s,r)$ are used to describe variations of the opening and the
traction along and across the interface.

\paragraph{Choice of the auxiliary cohesive-leg length \texorpdfstring{$\delta$}{delta}.}
The cohesive-leg length $\delta$ defines the available extent over which the
finite fracture-process zone may develop along the trial continuation.
Within the present formulation, $\delta$ is not a material parameter but an
auxiliary numerical length introduced to regularize the singular elastic
near-tip field by replacing it with a finite cohesive zone. For this
regularization to be meaningful, $\delta$ must be sufficiently long to
contain the full active cohesive zone. Choosing $\delta$ too short
artificially truncates the process zone and may leave a residual stress
concentration at the cohesive tip. Once the active cohesive part remains well
contained within the trial segment, however, $\delta$ no longer plays a
constitutive role and the predicted response becomes insensitive to its
precise value.

For trapezoidal cohesive laws, including the piecewise-smooth law adopted in
this work, a practical estimate of the required scale can be obtained from
the Dugdale model. In that idealization, the process-zone length
$\ell_p$ scales with the crack size $a$ as
\[
\ell_p = 2\pi a
\left(\frac{\sigma_{\mathrm{cr}}^{\mathrm{LEFM}}}{\sigma_n^{\max}}\right)^2.
\]
This estimate is used here as a practical guideline for selecting the
auxiliary cohesive length. In the baseline benchmark calculations reported in
Sections~\ref{sec:results} and~\ref{sec:LEFM-comparison}, a representative
fixed value $\delta/a=0.06$ is employed. In the sensitivity study of
Section~\ref{sec:scope_limitations}, by contrast, $\delta$ is chosen
adaptively from a calibrated LEFM-based predictor so that the active-zone
fraction $\beta/\delta$ remains controlled as the material parameters vary.
The role of $\delta$ should therefore be understood consistently throughout
the paper as that of an auxiliary regularization length, not as an intrinsic
fracture parameter.

Figure~\ref{fig:stmt:d} shows the mode-wise traction--separation laws used in
the cohesive formulation. In each pure mode, the traction increases from zero
to the cohesive strength and then decays smoothly to zero at complete
decohesion. The initial rising branch is not intended to represent material
hardening. Its purpose is to regularize the transition from the traction-free
crack faces to the active cohesive region. Without such a strengthening
branch, imposing a constant plateau traction immediately from $d_n=0$ would
create a jump in stress where the cohesive zone meets the traction-free part
of the crack, leading to an unphysical singularity at the process-zone tip.
The rising branch therefore ensures that the cohesive traction approaches its
maximum value smoothly, provides a finite-stress closure of the cohesive
zone, and allows the active cohesive length to emerge as part of the
solution rather than being prescribed in advance.

The main objective of the analysis is to determine the steady-state
direction of crack advance, or equivalently the kink angle, that develops in
response to the applied mixed-mode loading by evaluating competing trial
orientations of Leg~2 through the cohesive formulation. The subsequent
sections formulate the governing equations, define the mixed-mode cohesive
law quantitatively, and describe the finite-element construction used to
compute the corresponding critical states.


\section{Governing equations}
\label{sec:governing}

The displacement field
$\uu:\,\Omega\setminus(\Gamma_{\mathrm{cr}}\cup\Gamma_{\mathrm{cz}})\to\mathbb{R}^2$
is the primary unknown of the boundary-value problem.
The body is treated as a homogeneous, isotropic, linear-elastic solid in
plane strain, and equilibrium is assumed under quasi-static loading conditions.

\paragraph{Kinematics and constitutive law.}
The infinitesimal strain tensor and Cauchy stress tensor are related by the
linear elastic law
\begin{equation}
	\eps(\uu)
	=\tfrac12(\nabla\uu+\nabla\uu^\top),
	\qquad
	\sig=\four D:\eps(\uu),
	\label{eq:constitutive}
\end{equation}
where $\four D$ is the fourth-order isotropic elasticity tensor defined by
the Young’s modulus~$E$ and Poisson’s ratio~$\nu$.
Under plane-strain conditions the out-of-plane stress satisfies
$\sigma_{zz}=\nu(\sigma_{xx}+\sigma_{yy})$,
and the shear modulus is $G=E/[2(1+\nu)]$.


\paragraph{Equilibrium equations.}
In the absence of body forces, the equilibrium of stresses in the bulk
domain reads
\begin{equation}
	\nabla\!\cdot\!\sig=\vec 0
	\qquad
	\text{in } \Omega\setminus(\Gamma_{\mathrm{cr}}\cup\Gamma_{\mathrm{cz}}),
	\label{eq:equilibrium}
\end{equation}
where $\nabla\!\cdot(\cdot)$ denotes the row-wise divergence of a second-order
tensor.  Equation~\eqref{eq:equilibrium} governs the elastic field in all
regions of the body except along the discontinuity surfaces.

\paragraph{Boundary and interface conditions.}
Let $\vb n_b$ denote the outward unit normal to $\partial\Omega$.
Uniform tractions are applied on the horizontal edges (Fig.~\ref{fig:stmt:a}),
\[
\sig\,\vb n_b=
\begin{cases}
	+\sigma\,\vb e_y,& y=+B,\\
	-\sigma\,\vb e_y,& y=-B,
\end{cases}
\qquad\text{and}\qquad
\sig\,\vb n_b=\vb 0 \ \ \text{on } x=0 \text{ and } x=A,
\]
where the scalar $\sigma$ is the prescribed far-field load factor.
To suppress rigid-body motion we impose the minimal point constraints
\[
\uu(A,0)\!\cdot\!\vb e_y=0,\qquad
\uu(0,B)\!\cdot\!\vb e_x=0,\qquad
\uu(0,-B)\!\cdot\!\vb e_x=0,
\]
which fix one vertical and two horizontal degrees of freedom without otherwise
restricting the solution.

Across the faces of the physical crack $\Gamma_{\mathrm{cr}}$ the tractions
vanish:
\[
\sig\,\bn=\vb 0 \quad \text{on } \quad \Gamma_{\mathrm{cr}}^{\pm}.
\]
Along the cohesive segment $\Gamma_{\mathrm{cz}}$,
the displacement jump and cohesive traction satisfy the local interface relations
\begin{equation}
	\bd=\mat Q\,\jump{\uu},\qquad
	\bTcz=\bTcz(\bd) \in\mathbb{R}^2,
	\label{eq:cohesive_relations}
\end{equation}
\chg{R2C2}{Here, $\bTcz$ denotes the cohesive traction vector expressed as a nonlinear,
	history-independent function of the displacement jump $\bd$ across the interface,}
 $\mat Q=[\bt^\top;\bn^\top]$ is the local orthonormal frame attached to
Leg~2 (introduced in Section~\ref{sec:problem}),
$\bd=(d_t,d_n)$ are the local tangential sliding and normal opening displacements,
and $\bTcz=(\sigma_t,\sigma_n)$ are the corresponding cohesive tractions.
Traction continuity implies $\bTcz^{(+)}=-\,\bTcz^{(-)}$ along the interface,
and the functional dependence $\bTcz(\bd)$ is given by the potential-based
mixed-mode TSL in Section~\ref{sec:tsl}.


\paragraph{Weak form.}
Multiplying the equilibrium equation~\eqref{eq:equilibrium} by an admissible
virtual displacement field~$\delta\uu$ and integrating by parts yields the
variational (weak) form of equilibrium:
\begin{equation}
	\int_{\Omega\setminus(\Gamma_{\mathrm{cr}}\cup\Gamma_{\mathrm{cz}})}
	\eps(\delta\uu):\sig\,\dd\Omega
	-
	\int_{\Gamma_{\mathrm{cz}}}
	\delta\bd^\top\,\bTcz\,\dd s
	=
	\int_{\partial_t\Omega}
	\delta\uu^\top\,\sig\bn\,\dd s .
	\label{eq:weak_form}
\end{equation}
The first integral represents the bulk elastic contribution, while the second
integral accounts for the work of cohesive tractions over the displacement
jumps along $\Gamma_{\mathrm{cz}}$.  The right-hand side is the external work
of the applied tractions.  Equation~\eqref{eq:weak_form} constitutes the
starting point for the finite-element discretization developed in
Section~\ref{sec:construction}.

\paragraph{Remarks.}
The weak form separates the elastic and cohesive contributions explicitly,
making it possible to linearize the coupled system consistently with respect
to the interface degrees of freedom.  The cohesive-zone influence enters only
through the surface integral over $\Gamma_{\mathrm{cz}}$, which naturally
vanishes once the interface is fully degraded and the crack has propagated.
This formulation thus ensures a smooth transition from the intact to the
fully separated state without discontinuities in stiffness or traction fields.

\section{Cohesive interface and traction–separation law}
\label{sec:tsl}

The cohesive behavior along $\Gamma_{\mathrm{cz}}$ is governed by a
potential-based TSL. Local quantities are expressed
in the Leg--2 frame $\QQ=[\bt^\top;\bn^\top]$ introduced in
Section~\ref{sec:problem}. The face-to-face displacement jump resolves as
$\bd=(d_t,d_n)^\top=\QQ\,\jump{\uu}$, with normalized openings
$\bar d_t=|d_t|/d_t^{\max}$ and $\bar d_n=d_n/d_n^{\max}$.
Cohesive strengths are $\sigma_t^{\max}$ and $\sigma_n^{\max}$.
\chg{R2C3}{
	Together with the fracture energies $\phi_\alpha$
	($\alpha\in\{n,t\}$), the cohesive strengths $\sigma_\alpha^{\max}$
	constitute the independent fracture parameters.
	The critical openings $d_\alpha^{\max}$ are introduced as characteristic
	length scales and are determined from the fracture--energy condition
}
\begin{equation*}
	\mathcolor{Blue}{
		\phi_\alpha
		= \sigma_\alpha^{\max} d_\alpha^{\max} \, c_\alpha,
		\qquad
		c_\alpha = F_\alpha(1),
		\qquad
		F_\alpha(\xi)=\int_0^\xi t_\alpha(s)\,\mathrm{d}s .
	}
\end{equation*}
\chg{R2C3}{
	We further introduce the energetic prefactors
	\(
	\chi_\alpha := \sigma_\alpha^{\max} d_\alpha^{\max},
	\)
	which set the mode-wise energy scale of the cohesive potential.
	These quantities are not independent material parameters but follow
	directly from $(\sigma_\alpha^{\max},\phi_\alpha)$ and the chosen TSL shape.
}
\chg{R2C3}{
	Here $t_\alpha(\xi)$ denotes the normalized traction--separation law,
	$F_\alpha(\xi)$ the associated dimensionless potential, and
	$c_\alpha$ a dimensionless shape parameter that depends solely on the
	chosen TSL form and its parameters.
	For notational brevity, once the pure-mode traction--separation laws have
	been specified, we subsequently write $F(\cdot)$ and $t(\cdot)$ for the
	mode-dependent functions, with the active mode $\alpha$ uniquely determined
	by the argument and the associated parameter set
	$(\sigma_\alpha^{\max},\phi_\alpha,$ and the TSL shape parameters).
	No additional functional assumptions are implied by this notation.
}




\paragraph{Mode-wise (one-dimensional) law.}
The shape of the one-dimensional cohesive traction--separation response has a
direct influence on fracture behaviour. A cohesive law typically consists of a
rising branch, a peak traction, and a softening segment. Each part contributes
to the overall fracture process in a distinct manner: the initial strengthening
branch regularises the transition from contact to opening; the peak traction
governs crack initiation and the size of the FPZ; and the softening segment
controls the rate of energy dissipation and the interaction with the surrounding
elastic material. As emphasised by Shet and Chandra~\cite{ShetChandra2003}, the FPZ
embodies micromechanisms that are intrinsic to the material, so the
\emph{shape} of the traction--separation law, not merely the total energy or peak
traction, is a material parameter in its own right. To retain control over these
features while preserving smoothness and differentiability, we employ a
normalized traction $t(\xi)$ with a strengthening segment, a plateau, and a
softening branch.

\chg{R2C3}{
	The normalized traction--separation law $t$, its slope
	$g=t'$ and the associated potential $F$ are controlled by the parameters
	$0<a_1<a_2<1$, which define the end of the strengthening branch and the onset
	of softening, respectively. In applications, $a_1$ and $a_2$ (and hence the
	shape parameter $c$) may be chosen mode-wise; the subscript is suppressed here
	for brevity.
}

For the \emph{normal (tensile)} cohesive mode, a linear penalty is introduced in
compression to regularise contact, together with a $C^1$ contact--opening
transition enforced by $\kappa=2/a_1$. The resulting pure-mode law reads

\begin{multline}
	\label{eq:TSL-pure-mode}
	\mathcolor{Blue}{
	(F,t,g)(\xi)= \\
	\begin{cases}
		\big(\tfrac{\kappa}{2}\,\xi^2,\ \kappa\,\xi,\ \kappa\big),
		& \xi\le 0,\\[2pt]
		\big(\tfrac{\xi^2}{a_1}\bigl(1-\tfrac{\xi}{3a_1}\bigr),\
		\tfrac{\xi}{a_1}\bigl(2-\tfrac{\xi}{a_1}\bigr),\
		\tfrac{2}{a_1}\bigl(1-\tfrac{\xi}{a_1}\bigr)\big),
		& 0<\xi<a_1,\\[4pt]
		(\xi-\tfrac{a_1}{3},\ 1,\ 0),
		& a_1\le \xi<a_2,\\[4pt]
		\Big(
		c-\!\!\int_{\xi}^{1}\! t(s)\,\mathrm ds,\\
		\qquad\frac{(1+2\xi-3a_2)(1-\xi)^2}{(1-a_2)^3},\
		-\frac{6(1-\xi)(\xi-a_2)}{(1-a_2)^3}
		\Big),
		& a_2\le\xi\le 1,\\[8pt]
		(c,\ 0,\ 0),
		& \xi>1.
	\end{cases}
}
\end{multline}
\chg{R2C3}{Here $c=\tfrac{1}{6}(3-2a_1+3a_2)$ is the dimensionless area under the	normalized traction--separation law.
}
Choosing $\kappa=2/a_1$ yields $t(0^-)=t(0^+)=0$ and
$g(0^-)=g(0^+)=2/a_1$, i.e., a $C^1$ traction across the
contact--opening transition.



The dimensional pure-mode laws are written in terms of the normalized openings $\bar d_n$ and $\bar d_t$.
For the normal mode we use a tensile cohesive law with a contact penalty in
compression,
\[
\sigma_n(d_n)=
\begin{cases}
	\sigma_n^{\max}\,t(\bar d_n), & d_n\ge 0,\\[2pt]
	k_c\,d_n, & d_n<0,
\end{cases}
\qquad
k_c=\kappa\,\dfrac{\sigma_n^{\max}}{d_n^{\max}} .
\]
For the tangential mode we define the shear magnitude and the signed traction as
\[
|\sigma_t|(d_t)=\sigma_t^{\max}\,t(\bar d_t),
\qquad
\sigma_t(d_t)=|\sigma_t|(d_t)\,\operatorname{sgn}_0(d_t),
\]
where $\operatorname{sgn}_0(0)=0$.

\noindent Their (one-dimensional) tangents are
\[
\frac{\dd\sigma_n}{\dd d_n}=
\begin{cases}
	\displaystyle \dfrac{\sigma_n^{\max}}{d_n^{\max}}\,g(\bar d_n), & d_n>0,\\[6pt]
	k_c, & d_n<0,
\end{cases}
\qquad
\frac{\dd|\sigma_t|}{\dd |d_t|}=
\dfrac{\sigma_t^{\max}}{d_t^{\max}}\,g(\bar d_t).
\]

\chg{R2C3}{
\noindent The corresponding pure-mode potentials are taken as
\begin{equation*}
	\Psi_\alpha(d_\alpha)
	=\chi_\alpha\,F(\bar d_\alpha)
	=\frac{\phi_\alpha}{c_\alpha}\,F(\bar d_\alpha),
	\qquad \alpha\in\{n,t\}.
\end{equation*}
}

\paragraph{Potential-based mixed-mode coupling.}
The conservative mixed-mode cohesive response is formulated in terms of a single
scalar potential. \chg{R2C3}{ Building on the pure-mode energetic contributions
$\chi_\alpha F(\bar d_\alpha)$, the mixed-mode potential is defined as
\begin{equation}
	\label{eq:psi_mixed}
	\Psi(d_n,d_t)=
	\chi_n F(\bar d_n)
	+\chi_t F(\bar d_t)
	- r\,\sqrt{\chi_n\chi_t}\,
	F(\bar d_n)F(\bar d_t),
\end{equation}
The dimensionless parameter $r$ controls the strength of the mixed-mode
interaction. Similar multiplicative energetic couplings have been widely adopted in
potential-based mixed-mode cohesive formulations.
Admissibility of the cohesive law requires the multiplicative degradation
factors in the resulting tractions to remain non-negative over the admissible
range $0\le F(\xi)\le c_\alpha$. This yields the bound
\[
0\le r \le r^{\max}
=\min\!\Bigg(
\frac{1}{c_n}\sqrt{\frac{\chi_n}{\chi_t}},
\frac{1}{c_t}\sqrt{\frac{\chi_t}{\chi_n}}
\Bigg).
\]
}

\chg{R2C5}{
	The form \eqref{eq:psi_mixed} constitutes a minimal energetic extension of
	the pure-mode cohesive laws. It is derived from a single scalar potential,
	ensures variational consistency, and exactly recovers the uncoupled normal or
	tangential response when the complementary separation vanishes.
}

The cohesive tractions follow from
$\bTcz=\partial\Psi/\partial\bd
=(\sigma_t,\sigma_n)^\top$ as
\chg{R2C3}{
\begin{align}
	\label{eq:tractions_mixed}
	\begin{bmatrix}\sigma_t\\[2pt]\sigma_n\end{bmatrix}
	&=
	\begin{bmatrix}
		\sigma_t^{\max}
		\!\left(
		1-r\sqrt{\dfrac{\chi_n}{\chi_t}}\,F(\bar d_n)
		\right)
		t(\bar d_t)\,\operatorname{sgn}_0(d_t)
		\\[10pt]
		\sigma_n^{\max}
		\!\left(
		1-r\sqrt{\dfrac{\chi_t}{\chi_n}}\,F(\bar d_t)
		\right)
		t(\bar d_n)
	\end{bmatrix}.
\end{align}
}
\paragraph{Consistent tangent (local stiffness).}
Because $\Psi$ is a scalar potential, the consistent local tangent is the
symmetric Hessian
$\mat K_{\rm loc}=\partial\bTcz/\partial\bd$.
It admits the exact symmetric diagonal factorization
\begin{equation}
	\label{eq:Kloc-factor}
	\mat K_{\rm loc}=\mat S\,\mat M(\bar d_n,\bar d_t)\,\mat S,
	\qquad
	\mat S=\operatorname{diag}\!\left(
	\sqrt{\tfrac{\sigma_t^{\max}}{d_t^{\max}}},\
	\sqrt{\tfrac{\sigma_n^{\max}}{d_n^{\max}}}
	\right),
\end{equation}
where the dimensionless symmetric \emph{core} depends only on the normalized
openings:
\chg{R2C3}{
\begin{equation}
	\label{eq:Kcore}
	\mat M(\bar d_n,\bar d_t)=
	\begin{bmatrix}
		\Big[1-r\sqrt{\tfrac{\chi_n}{\chi_t}}\,F(\bar d_n)\Big]\,g(\bar d_t)
		&
		-\,r\,t(\bar d_t)\,t(\bar d_n)
		\\[6pt]
		-\,r\,t(\bar d_t)\,t(\bar d_n)
		&
		\Big[1-r\sqrt{\tfrac{\chi_t}{\chi_n}}\,F(\bar d_t)\Big]\,g(\bar d_n)
	\end{bmatrix}.
\end{equation}
}
This structure is key: the tangent is symmetric by construction (as the Hessian
of a potential), which improves robustness and supports quadratic convergence in
the Newton iterations used later.



\paragraph{Coupled response illustration.}
The coupling effects introduced by Eq.~\eqref{eq:psi_mixed} are illustrated in
Figs.~\ref{fig:tsl_2d} and~\ref{fig:tsl_3d}.
Figure~\ref{fig:tsl_2d} shows the normalized normal and tangential tractions
for several fixed values of the companion separation. As the tangential
separation $\bar d_t$ increases, the normal traction $\bar\sigma_n$ is reduced,
and vice versa, indicating gradual degradation of the complementary mode.
Figure~\ref{fig:tsl_3d} presents the mixed-mode cohesive potential in normalized
form,
\[
\hat\Psi(\bar d_n,\bar d_t)=\Psi(\bar d_n,\bar d_t)/\sqrt{\chi_n\chi_t},
\]
together with the corresponding tractions and the components of the local
tangent matrix~$\mat K_{\rm loc}$. This normalization is used for visualization only and does not enter the
governing equations or the numerical implementation.
The coupling between the normal and tangential directions appears as the
nonzero off-diagonal term~$K_{nt}$.




\begin{figure}[t]
	\centering
	\includegraphics[width=0.98\textwidth]{Fig2New}
	\caption{Mode coupling in the potential-based TSL.
		Left: variation of normalized normal traction~$\bar\sigma_n$ with~$\bar d_n$
		for fixed~$\bar d_t$.
		Right: variation of normalized tangential traction~$\bar\sigma_t$
		with~$\bar d_t$ for fixed~$\bar d_n$.
		Increasing separation in one mode reduces the cohesive resistance in the
		other.}
	\label{fig:tsl_2d}
\end{figure}

\begin{figure}[t]
	\centering
	\includegraphics[width=\textwidth]{Fig3New}
	\caption{Potential-based mixed-mode traction–separation law.
		Top: cohesive potential~$\Psi(\bar d_n,\bar d_t)$ and normalized tractions.
		Bottom: components of the local tangent stiffness matrix~$\mat K_{\rm loc}$,
		showing both the diagonal terms and the coupling~$K_{nt}$.}
	\label{fig:tsl_3d}
\end{figure}

\paragraph{Parameters used in Figs.~\ref{fig:tsl_2d}–\ref{fig:tsl_3d}.}
We take the \emph{dimensionless} TSL shape and coupling parameters
\[
a_1=0.2,\qquad a_{2n}=0.6,\qquad a_{2t}=0.3,\qquad r=0.40,
\]
and the following \emph{dimensional} material/TSL scales:
\[
\sigma_n^{\max}=40~\text{MPa},\quad \sigma_t^{\max}=16~\text{MPa},\quad
\phi_n = 800~\text{J\,m}^{-2},\quad \phi_t = 1200~\text{J\,m}^{-2}.
\]

For the normalized law \eqref{eq:TSL-pure-mode}, the dimensionless areas under
$t(\xi)$ are
\[
c_n=\frac{1}{6}\bigl(3-2a_1+3a_{2n}\bigr)=0.7333,
\qquad
c_t=\frac{1}{6}\bigl(3-2a_1+3a_{2t}\bigr)=0.5833.
\]

The critical openings follow from the fracture--energy condition
$d_\alpha^{\max}= \phi_\alpha/(\sigma_\alpha^{\max}c_\alpha)$, yielding
\[
d_n^{\max}\approx 27.3~\mu\text{m},
\qquad
d_t^{\max}\approx 129~\mu\text{m}.
\]

The mixed-mode coupling parameter is admissible under the bound
$r\le r^{\max}\approx 1.25$; thus $r=0.40$ is admissible.

The $C^1$ compression--opening switch on the normal branch uses the
(dimensionless) penalty
\(
\kappa={2}/{a_1}=10,
\)
which corresponds to the dimensional contact slope
$k_c=\kappa\,\sigma_n^{\max}/d_n^{\max}$ (units: Pa\,m$^{-1}$), applied only for $d_n<0$.
All figure fields are displayed in normalized form, i.e., tractions scaled by
$(\sigma_n^{\max},\sigma_t^{\max})$ and openings by $(d_n^{\max},d_t^{\max})$,
while the parameters above retain their physical units.





\paragraph{Remark on the choice of $a_{1}$.}
The parameter $a_{1}$ controls the end of the strengthening branch of the
traction--separation law and therefore sets the initial crack opening over
which cohesive tractions rise to their peak value.
In Section~\ref{sec:tsl}, a comparatively large value was used for clarity of
presentation, but such values are not representative of real materials: if
$a_{1}$ is too large, the apparent process-zone opening becomes unrealistically
extended.
In practice, $a_{1}$ should be chosen as small as possible while still
ensuring rapid and stable convergence of the nonlinear solver.
\chg{R1C2}{
    For the range of values considered in this work, no locking effects or spurious
    stiffening were observed as $a_{1}$ was reduced, including under global mesh
    refinement; the primary effect of decreasing $a_{1}$ is an increase in the
    number of Newton iterations required to reach convergence.
}
The high-resolution studies in Section~\ref{sec:results} show that
$a_{1}\lesssim 0.01$ provides accurate traction reconstruction and fast
convergence, whereas larger values artificially enlarge the crack opening
across the strengthening interval (the $(\beta,\delta)$ portion of the cohesive
zone, see Fig.~\ref{fig:stmt:c}).
\chg{R1C2}{
    The potential-based formulation employed here yields a symmetric and consistent
    algorithmic tangent, which stabilizes the Newton process and reduces sensitivity
    to the specific choice of $a_{1}$ compared to non-potential-based cohesive laws.
}



\paragraph{Discussion of coupling effects.}
The potential~\eqref{eq:psi_mixed} ensures energy consistency and a symmetric
tangent matrix; however, the coupling it introduces also generates certain
nonphysical features.
When the tangential separation becomes critical while the normal opening
remains small, the model still produces finite normal cohesive tractions.
Conversely, when the interface opens completely in the normal direction,
tangential tractions do not vanish entirely.
This residual stress field is a known limitation of potential-based TSLs
because energy coupling prevents either traction component from reaching zero
independently.  In practice, the effect is small for moderate~$r$
($r\le0.3$), but it should be considered when interpreting mixed-mode
crack-path predictions.

\paragraph{Significance of the potential.}
The cohesive potential provides a single, smooth energetic surface from which
all tractions and tangent moduli follow by differentiation.
Its existence guarantees that the interface response is conservative and
thermodynamically consistent, and that the linearization used in
Section~\ref{sec:construction} yields a symmetric tangent operator.
This symmetry improves numerical stability and ensures quadratic convergence
of Newton iterations while maintaining a physically consistent distribution of
cohesive energy among modes.

\chg{R1C6}{
	\paragraph{Relation to stabilized and energetically consistent CZM formulations.}
	Numerical robustness and energetic consistency of cohesive-zone formulations
	have been extensively studied in the literature.
	Stabilized interface-element approaches have been proposed, for example by
	Ghosh et al.~\cite{Ghosh2019StabilizedCZM,Ghosh2021RobustnessCZM}, to address
	convergence difficulties arising in delamination analyses when very stiff,
	penalty-like cohesive laws are employed.
	From a complementary perspective, Dimitri et al.~\cite{Dimitri2014ConsistencyCZM}
	emphasized the importance of thermodynamic consistency and potential-based
	formulations in mixed-mode cohesive models to avoid non-conservative or
	path-dependent responses.
	The present work aligns with this latter viewpoint: robustness is achieved
	intrinsically through a unified potential-based cohesive law with a symmetric
	Hessian tangent, rather than through additional stabilization terms or extreme
	penalty stiffness.
	This guarantees consistent energetic coupling between modes and stable Newton
	iterations while retaining a clear physical interpretation of the cohesive
	parameters.
}


\section{Cohesive interface discretization and consistent linearization}
\label{sec:construction}

We discretize the cohesive interface $\Gamma_{\mathrm{cz}}$ by keeping the two
faces of Leg~2 kinematically distinct and pairing coincident node triplets
along a thin channel (Section~\ref{sec:mesh}).  The element operators below are
given for a \emph{quadratic} interface segment $e$ with coincident
\emph{top} nodes $(i^+,j^+,k^+)$ and \emph{bottom} nodes $(i^-,j^-,k^-)$.
The segment is oriented from $i$ (mouth) to $k$ (tip). A local orthonormal
frame $\QQ=[\bt^\top;\bn^\top]$ is attached tangentially (mouth$\to$tip) and
normally (lower$\to$upper), see Fig.~\ref{fig:coh-tri}. Within each interface segment, $\QQ$ is taken constant and evaluated from the segment chord (mouth$\to$tip).


\begin{figure}[!h]
	\centering
	\includegraphics[width=.58\textwidth]{Fig4}
	\caption{Quadratic interface segment with coincident node triplets
		$(i^\pm,j^\pm,k^\pm)$.  The top edge is $(i^+,j^+,k^+)$; the bottom edge is
		numbered CCW as $(k^-,j^-,i^-)$.
		Local frame $\QQ=[\bt^\top;\bn^\top]$ is attached to the segment.}
	\label{fig:coh-tri}
\end{figure}

\paragraph{Unknown ordering, shape rows, and edge permutation.}
We stack unknowns as \emph{top then bottom}, with the bottom group reversed
to follow the lower element’s CCW edge order:
\begin{align}
	\vb U_e&=\big[\ \vb U_e^{+}\ \ \vb U_e^{-}\ \big]^\top\in\mathbb{R}^{12}, \label{eq:seg-Ue}\\
	\vb U_e^{+}&=\big[\,u_{i^+x}\ u_{i^+y}\ u_{j^+x}\ u_{j^+y}\ u_{k^+x}\ u_{k^+y}\,\big]^\top,\notag\\
	\vb U_e^{-}&=\big[\,u_{k^-x}\ u_{k^-y}\ u_{j^-x}\ u_{j^-y}\ u_{i^-x}\ u_{i^-y}\,\big]^\top.\notag
\end{align}
Let
\[
N=\begin{bmatrix}N(\xi_i)\\ N(\xi_j)\\ N(\xi_k)\end{bmatrix},\qquad
N(\xi_\ell)=\big[N_i(\xi_\ell)\ \ N_j(\xi_\ell)\ \ N_k(\xi_\ell)\big],\quad
\ell\in\{i,j,k\},
\]
be the stacked scalar quadratic shape rows, and define the node-order
reversal
\begin{equation*}
	\mat P_e=\begin{bmatrix}0&0&1\\0&1&0\\1&0&0\end{bmatrix},
	\qquad (i,j,k)\mapsto(k,j,i).
	\label{eq:seg-Pi}
\end{equation*}

\paragraph{Explicit jump operator.}
Stack the local openings at the three segment nodes as
\(\bd_e=[\,d_{t,i}\ d_{n,i}\ d_{t,j}\ d_{n,j}\ d_{t,k}\ d_{n,k}\,]^\top\in\mathbb{R}^6\).
With the CCW ordering in Eq.~\eqref{eq:seg-Ue}, the nodal jump map is
\begin{align}
	&\bd_e=\mat B_e\,\vb U_e,\qquad
	\mat B_e=\big[\,\mat B_e^{+}\ \ \mat B_e^{-}\,\big], \label{eq:seg-Be}\\
	&\mat B_e^{+}=(\mat I_3\otimes\QQ)\,(N\otimes \mat I_2),\qquad
	\mat B_e^{-}=-(\mat I_3\otimes\QQ)\,\big((N\mat P_e)\otimes \mat I_2\big), \notag
\end{align}
with $B_e$ extracting the opening displacement across the interface. In block form (each entry $2{\times}2$),
\[
\mat B_e=
\begin{bmatrix}
	N_i(\xi_i)\QQ & N_j(\xi_i)\QQ & N_k(\xi_i)\QQ & -N_k(\xi_i)\QQ & -N_j(\xi_i)\QQ & -N_i(\xi_i)\QQ\\
	N_i(\xi_j)\QQ & N_j(\xi_j)\QQ & N_k(\xi_j)\QQ & -N_k(\xi_j)\QQ & -N_j(\xi_j)\QQ & -N_i(\xi_j)\QQ\\
	N_i(\xi_k)\QQ & N_j(\xi_k)\QQ & N_k(\xi_k)\QQ & -N_k(\xi_k)\QQ & -N_j(\xi_k)\QQ & -N_i(\xi_k)\QQ
\end{bmatrix}.
\]
\emph{Lobatto nodal integration used in the implementation}  (evaluation at $\xi\in\{\xi_i,\xi_j,\xi_k\}$, hence $N=\mat I_3$):
\[
\mat B_e=
\begin{bmatrix}
	\QQ & \bm 0 & \bm 0 & \bm 0 & \bm 0 & -\QQ\\
	\bm 0 & \QQ & \bm 0 & \bm 0 & -\QQ & \bm 0\\
	\bm 0 & \bm 0 & \QQ & -\QQ & \bm 0 & \bm 0
\end{bmatrix}.
\]

\paragraph{Line integration, rotation, and action–reaction.}
Let $\xi\in[0,1]$ denote the parent coordinate along the segment and $l_e$ its physical length.
With the segment’s consistent line “mass’’
\begin{equation}
	\mat R=l_e\!\int_0^1\!N(\xi)^\top N(\xi)\,\dd\xi
	=\frac{l_e}{30}
	\begin{bmatrix}
		4&2&-1\\[2pt]
		2&16&2\\[2pt]
		-1&2&4
	\end{bmatrix},
\end{equation}
the \emph{top-face} nodal resultants are
\[
\vb t_e^{+}=(\mat R\otimes \QQ^\top)\,\vb s_e\in\mathbb{R}^{6},
\]
\chg{R2C7}{where
	$\vb s_e=\big[\,\sigma_{t,i}\ \sigma_{n,i}\ \sigma_{t,j}\ \sigma_{n,j}\ \sigma_{t,k}\ \sigma_{n,k}\,\big]^\top$
	stacks the six local cohesive tractions at the three evaluation stations (chosen as the Lobatto nodes in the implementation).
	These tractions are obtained from the corresponding local displacement jumps
	via the nonlinear traction--separation mapping
	$\mathcal T:\mathbb{R}^6\to\mathbb{R}^6$, i.e. 
	\begin{equation}
	\vb s_e=\mathcal T(\bd_e).
		\end{equation}
	}
At each node $\ell\in\{i,j,k\}$ we use the pointwise traction map (see Eq.~\eqref{eq:tractions_mixed} from
		Section~\ref{sec:tsl}):
		\begin{align}
			\label{eq:seg-TSL-point}
			&\begin{bmatrix}\sigma_{t,\ell}\\[2pt]\sigma_{n,\ell}\end{bmatrix}
			= \bTcz \big(\bar d_{n,\ell},\,\bar d_{t,\ell};\,s_{t,\ell}\big),\\
			&\bar d_{t,\ell}=\frac{|d_{t,\ell}|}{d_t^{\max}},\quad
			\bar d_{n,\ell}=\frac{d_{n,\ell}}{d_n^{\max}},\qquad
		s_{t,\ell}=\operatorname{sgn}_0(d_{t,\ell}).
		\notag
		\end{align}

The \emph{bottom-face} resultants, equal and opposite but stored in the
flipped order $(k^-,j^-,i^-)$, are
\[
\vb t_e^{-}=-\,(\mat P_e\otimes \mat I_2)\,\vb t_e^{+}.
\]
Stacking both faces yields the $12$-component element residual; using the
traction–separation mapping $\mathcal T$ and the jump $\bd_e=\mat B_e\vb U_e$:
\begin{equation}
	\vb f_e
	=
	\underbrace{\begin{bmatrix} \mat I_6\\[2pt]-\,(\mat P_e\otimes \mat I_2)\end{bmatrix}(\mat R\otimes \QQ^\top)}_{\displaystyle \mat A_e}
	\;
	\underbrace{\mathcal T\!\big(\mat B_e\,\vb U_e\big)}_{\displaystyle \vb s_e}
	\ \in\ \mathbb{R}^{12}.
	\label{eq:seg-res}
\end{equation}
Here $\mat A_e$ collects the line weighting $\mat R$, rotation $\QQ^\top$, and the action--reaction stacking of the two faces.



\paragraph{Consistent linearization (Jacobian).}
From the residual composition in Eq.~\eqref{eq:seg-res}, the chain rule yields
the cohesive element tangent
\begin{equation}
	\mat K_e
	=\frac{\partial \vb f_e}{\partial \vb U_e}
	=
	\underbrace{\frac{\partial \vb f_e}{\partial \vb s_e}}_{\displaystyle \mat A_e}
	\underbrace{\frac{\partial \vb s_e}{\partial \bd_e}}_{\displaystyle \mat K_{\rm loc}}
	\underbrace{\frac{\partial (\mat B_e \vb U_e)}{\partial \vb U_e}}_{\displaystyle \mat B_e}
	\;=\;
	\mat A_e\,\mat K_{\rm loc}\,\mat B_e
	\;\in\;\mathbb{R}^{12\times12}.
	\label{eq:seg-J}
\end{equation}
The block-diagonal local tangent collects the pointwise contributions at the
quadratic cohesive stations,
\[
\mat K_{\rm loc}
=
\operatorname{blkdiag}(\mat K_i,\mat K_j,\mat K_k),
\]
The pointwise tangent \(\mat K_\ell\) follows from the
potential-based cohesive law described in Section~\ref{sec:tsl}.
The use of a scalar cohesive potential guarantees that the local core tangent
is symmetric, and this symmetry is preserved in the assembled element Jacobian
away from the shear sign switch at \(d_t=0\).



\paragraph{Global assembly.}
Assembling the element vectors $\vb f_e$ and Jacobians $\mat K_e$
over all interface segments yields the global cohesive residual and tangent
\[
\vb F_{\mathrm{coh}}(\vb u)\in\mathbb{R}^{n_{\mathrm{dof}}},
\qquad
\frac{\partial \vb F_{\mathrm{coh}}}{\partial \vb u}
\in\mathbb{R}^{n_{\mathrm{dof}}\times n_{\mathrm{dof}}}.
\]
The bulk stiffness $\mat K$ enters additively in the equilibrium block.

\paragraph{Augmented equilibrium and Jacobian.}
We solve for the displacement vector
$\vb u\in\mathbb{R}^{n_{\mathrm{dof}}}$ and the scalar load amplitude
$\sigma\in\mathbb{R}$ from the augmented system
\[
\vb F_{\mathrm{aug}}(\vb u,\sigma)=
\begin{bmatrix}
	\vb R(\vb u,\sigma)\\[4pt] g(\vb u)
\end{bmatrix}
=\chg{R2C8}{
	\begin{bmatrix}
		\{\mat K\,\vb u - \sigma\,\vb F_{\mathrm{ext}} - \vb F_{\mathrm{coh}}(\vb u)\\[4pt]
		d_n^{\text{tip}} - d_n^{\max}
	\end{bmatrix}=\vb 0,}
\]
where $\vb F_{\mathrm{ext}}\in\mathbb{R}^{n_{\mathrm{dof}}}$ is the unit load pattern and
$g:\mathbb{R}^{n_{\mathrm{dof}}}\!\to\mathbb{R}$ extracts the normal opening
$d_n^{\text{tip}}$ at the cohesive tip (with four nonzero entries corresponding
to the two tip nodes).

The Newton Jacobian has the block form
\[
\mat J_{\mathrm{aug}}=
\frac{\partial \vb F_{\mathrm{aug}}}{\partial(\vb u,\sigma)}=
\begin{bmatrix}
	\displaystyle \mat K - \frac{\partial \vb F_{\mathrm{coh}}}{\partial \vb u}
	& \displaystyle -\,\vb F_{\mathrm{ext}}\\[10pt]
	\displaystyle \frac{\partial g}{\partial \vb u} & 0
\end{bmatrix}
\in \mathbb{R}^{(n_{\mathrm{dof}}+1)\times(n_{\mathrm{dof}}+1)}.
\]
Here $\dfrac{\partial g}{\partial \vb u}\in\mathbb{R}^{1\times n_{\mathrm{dof}}}$
has four nonzeros,
\[
\frac{\partial g}{\partial \vb u} =
\begin{bmatrix}
	\cdots\; \bn^{\!\top}\; \cdots\; -\bn^{\!\top}\; \cdots
\end{bmatrix},
\qquad
\bn=\text{(local normal)}=\text{row }2\text{ of }\QQ\in\mathbb{R}^{2\times 2}.
\]

\paragraph{Newton update.}
At iteration $k$, we solve
\[
\mat J_{\mathrm{aug}}(\vb u^{(k)},\sigma^{(k)})
\begin{bmatrix}\Delta\vb u\\ \Delta\sigma\end{bmatrix}
= -\,\vb F_{\mathrm{aug}}(\vb u^{(k)},\sigma^{(k)}),
\]
and update
$\vb u^{(k+1)}=\vb u^{(k)}+\Delta\vb u$,
$\sigma^{(k+1)}=\sigma^{(k)}+\Delta\sigma$,
until the relative residual norm and the cohesive traction norm fall below the
prescribed tolerances.


\paragraph{Concluding notes.}
The interface formulation separates, by design, (i) the kinematic jump map
$\mat B_e$, (ii) the local cohesive evaluation $\mathcal T$, and (iii) the
symmetric, potential-consistent tangent $\mat K_{\rm loc}=\partial\bTcz/\partial\bd$,
assembled via $\mat K_e=\mat A_e\,\mat K_{\rm loc}\,\mat B_e$.
This modular structure makes it straightforward to swap traction–separation
variants or add rate effects without altering the FE scaffolding, while the
Hessian form of $\mat K_{\rm loc}$ guarantees symmetry and robust Newton
convergence.
With the element operators defined, we now specify how the two-leg path is
embedded into the domain mesh—channel width, seeding of Leg~2, and graded
sizes—in Section~\ref{sec:mesh}.

\section{Mesh creation}
\label{sec:mesh}

The mesh must resolve the narrow cohesive segment (Leg~2) while keeping the
global model compact. We generate a thin channel solely to drive \emph{triangular} (T3) refinement
along the two-leg path; after meshing, we convert to quadratic (T6) elements
and collapse the channel to obtain a zero-thickness crack with collocated
top/bottom nodes.


\paragraph{Channel geometry and face topology (for T3 meshing).}
A narrow channel is built around $\Gamma_{\mathrm{cr}}\cup\Gamma_{\mathrm{cz}}$
solely for T3 mesh generation. Its half-width is tied to the interface
resolution,
\[
w_c=\frac{h_{\min}}{16},
\]
chosen empirically to balance refinement enforcement and element quality.
 Along Leg~2 the two inner faces are \emph{not merged} during
T3 generation; they are meshed as coincident edges with separate node sets so
that the displacement jump $\bd=\QQ\,\jump{\uu}$ is well defined at nodes.

\paragraph{Target sizing and uniform seeding of Leg~2.}
Leg~2 is uniformly seeded into $n_{\mathrm{coh}}$ stations, which fixes the
near-interface target size
\(
h_{\min}={\delta}/{n_{\mathrm{coh}}}.
\)
The physical crack (Leg~1) is seeded consistently with the adjacent bulk
resolution to make the transition into the cohesive segment smooth.

\paragraph{Graded size field and far field cap.}
Away from the channel, the target element size grows geometrically with factor
$h_{\mathrm{grad}}>1$ up to a far-field cap $h_{\max}$. This produces a smooth
gradation from fine elements near Leg~2 to coarser elements in the outer
regions, avoiding stiffness jumps and keeping the system size moderate.

\paragraph{T3\,$\to$\,T6 conversion and channel collapse.}
After T3 generation, all bulk triangles are upgraded to T6. The narrow channel
is then \emph{collapsed} to the polyline midline so that the two faces along
Leg~2 coincide geometrically while remaining kinematically distinct via
coincident node pairs (top/bottom). Interface segments use quadratic elements with Lobatto nodal evaluation
($N=\mat I_3$), matching the operators of Section~\ref{sec:construction}
so that $\bd_e=\mat B_e\vb U_e$ and tractions are evaluated pointwise at
segment nodes.


\paragraph{Quality and topology checks.}
The band is refined to maintain acceptable aspect ratios and to avoid slivers
near the knee $V_1$. Face orientation is enforced so the local frame
$\QQ=[\bt^\top;\bn^\top]$ points from lower to upper consistently
(mouth$\to$tip along $\bt$). Coincident node pairing along Leg~2 is verified
prior to assembly to guarantee one-to-one cohesive linkage. 

\paragraph{Verification and default parameters.}
Mesh independence is assessed by halving $h_{\min}$ (and, if needed, adjusting
$h_{\max}$ proportionally) while keeping $h_{\mathrm{grad}}$ fixed. The
critical load $\sigma_{\mathrm{cr}}$ and kink angle $\theta_2$ are deemed
converged when relative changes between successive refinements are below $1\%$.
We use $w_c=h_{\min}/16$ with
$h_{\min}=\delta/n_{\mathrm{coh}}$.

\paragraph{Illustration.}
\chg{R2C9}{
	Figure~\ref{fig:mesh:a} shows the pre-collapse thin channel constructed around
	the prescribed two-leg polyline, which enforces local mesh refinement along the
	intended cohesive path.
}
Figure~\ref{fig:mesh:b} shows the post-collapse zero-thickness crack with collocated top/bottom nodes and quadratic edges.
In the example we use $n_{\text{coh}}=20$ (so $h_{\min}=\delta/n_{\text{coh}}$), $\delta=0.06\,a$, $h_{\mathrm{grad}}=1.2$,
and $h_{\max}=B/5$ ($B$ is the plate half-height).

\begin{figure}[h!]
	\centering
	\begin{subfigure}[t]{0.48\textwidth}
		\centering
		\includegraphics[width=\linewidth]{Fig5a.jpg}
		\caption{Pre-collapse thin channel (refinement band).}
		\label{fig:mesh:a}
	\end{subfigure}\hfill
	\begin{subfigure}[t]{0.48\textwidth}
		\centering
		\includegraphics[width=\linewidth]{Fig5b.jpg}
		\caption{Post-collapse zero-thickness crack (T6).}
		\label{fig:mesh:b}
	\end{subfigure}
	\caption{Mesh near the two-leg crack before and after channel collapse.}
	\label{fig:mesh}
\end{figure}

\chg{R1C1}{
    \paragraph{Remark on channel collapse versus split-node insertion.}
    The present channel-collapse strategy is closely related in spirit to standard
    split-node procedures for introducing zero-thickness cohesive interfaces, but
    differs in the stage at which mesh resolution is enforced.
    In classical split-node approaches, the interface is created \emph{after}
    meshing by duplicating nodes and modifying element connectivity along a
    prescribed path; while effective for simple, straight cracks, this does not
    by itself control the resolution of bulk fields normal to the interface and
    often requires additional local remeshing.
    In contrast, the channel-collapse approach constructs the mesh \emph{a priori}
    with a thin refinement band aligned with the intended cohesive leg, so that the
    interface discretization and the surrounding bulk gradation are governed
    simultaneously through the size field parameters
    $(h_{\min},h_{\mathrm{grad}},h_{\max})$.
    The subsequent collapse step produces a zero-thickness interface with
    coincident node pairs and ordered topology without explicit mesh surgery.
    This entails a modest preprocessing effort to define the channel geometry, but
    provides improved robustness, reproducibility, and resolution control in
    systematic crack-kinking studies involving multiple trial paths.
}


\paragraph{Transition to results.}
With the interface operators (Section~\ref{sec:construction}) and the graded,
collapsed mesh defined here, the next section reports verification, mesh
convergence, and kink-angle selection under mixed-mode loading.

\section{Results}
\label{sec:results}

For each trial direction $\theta_2$ of Leg~2 we solve the augmented system
to the \emph{critical state}, defined by the Leg--2 mouth reaching the normal
critical separation $d_n^{\max}$. Newton iterations employ the consistent
cohesive tangent from Section~\ref{sec:tsl}. Interface fields are post-processed
in the local frame $(\bt,\bn)$ attached to Leg~2. Field plots use normalized variables
$\bar\sigma_{n,t}=\sigma_{n,t}/\sigma_{n,t}^{\max}$ and
$(\bar d_n,\bar d_t)=(d_n/d_n^{\max},\,|d_t|/d_t^{\max})$.
The \emph{critical load} $\sigma_{\mathrm{cr}}$ reported in tables and sweeps
is dimensional (MPa).

\paragraph{Parameters.}
The angle of the physical crack is $\theta_1=-20^\circ$.
Plane strain is assumed with $E=4$\,GPa and $\nu=0.3$.
Cohesive strengths
$(\sigma_n^{\max},\sigma_t^{\max})=(40,\ 15.4)$\,MPa
(from $\sigma_n^{\max}=E/100$ and $\sigma_t^{\max}=G/100$ with
$G=E/[2(1+\nu)]$).
The prescribed fracture energies are
$(\phi_n,\phi_t)=(800,\ 1200)$\,J\,m$^{-2}$,
with coupling parameter $r=0.4$.
One-dimensional TSL shape parameters are $a_1=0.002$,
$a_{2n}=0.9$, and $a_{2t}=0.5$.
\chg{R2C3}{Using the corresponding shape parameter $c$ for each mode,
the resulting critical openings are
$d_n^{\max}\approx 23.6\,\mu\text{m}$ and
$d_t^{\max}\approx104\,\mu\text{m}$.
}

\subsection{Convergence}
\label{sec:conv}

Convergence is assessed with respect to the interface resolution
$n_{\mathrm{coh}}$, such that the minimum characteristic length along the
cohesive leg is $h_{\min}=\delta/n_{\mathrm{coh}}$.
Two complementary studies are performed:
(i) refinement of the cohesive discretization at fixed far-field truncation,
and (ii) refinement of the far-field cap at fixed interface resolution.
All results reported in this section correspond to a representative mixed-mode
configuration with $\theta_2=20^\circ$.

\paragraph{Interface refinement.}
At fixed far-field cap $B/h_{\max}=80$, the cohesive interface is refined as
$n_{\mathrm{coh}}=\{10,20,40,80\}$.
\chg{R1C3}{ Table~\ref{tab:conv_far80} reports the number of degrees of freedom, the critical
load $\sigma_{\mathrm{cr}}$, and the ratio between mesh-generation time and
nonlinear solution time, $t_{\mathrm{mesh}}/t_{\mathrm{solve}}$.
This ratio provides a direct measure of the relative cost of remeshing compared
to solving the nonlinear equilibrium problem.}

With increasing $n_{\mathrm{coh}}$, the critical load rapidly stabilises.
Applying a three-level Richardson/Aitken extrapolation to the finest three
levels yields an asymptotic limit
$\sigma_\infty = 5.875349$~MPa and a finest-mesh relative error
\[
\epsilon
=
\frac{\big|\sigma_{\mathrm{cr}}^{\text{(finest)}}-\sigma_\infty\big|}
{\big|\sigma_\infty\big|}
\simeq 1.8\times10^{-6},
\]
indicating that the calculation has entered the asymptotic regime with respect
to interface refinement.
\chg{R1C3}{Over the full range of resolutions considered, the cost of remeshing remains
comparable to, but does not exceed, the cost of the nonlinear solve.}

\begin{table}[h!]
	\centering
	\caption{Convergence with interface refinement at fixed far-field cap
		$B/h_{\max}=80$ and $\theta_2=20^\circ$.}
	\label{tab:conv_far80}
	\begin{tabular}{rcccc}
		\toprule
		$n_{\mathrm{coh}}$ & DOFs & $\sigma_{\mathrm{cr}}$ [MPa]
		& $t_{\mathrm{mesh}}/t_{\mathrm{solve}}$ & iters \\
		\midrule
		10 & 120{,}962 & 5.81448 & 1.09 & 17 \\
		20 & 123{,}446 & 5.90673 & 1.15 & 16 \\
		40 & 127{,}162 & 5.87593 & 1.05 & 17 \\
		80 & 133{,}184 & 5.87536 & 1.21 & 17 \\
		\bottomrule
	\end{tabular}
\end{table}

\paragraph{Far-field truncation.}
To assess sensitivity to the outer boundary, a second sweep is performed by
varying the far-field cap
$B/h_{\max}\in\{10,20,40,80\}$ while keeping the cohesive resolution fixed at
$n_{\mathrm{coh}}=80$.
The results are summarised in Table~\ref{tab:farfield_sweep}.

As the far-field cap is increased, the critical load decreases monotonically
and converges towards the same asymptotic value obtained in the interface
refinement study.
The difference between $B/h_{\max}=40$ and $80$ is below $1.3\times10^{-4}$~MPa,
confirming that far-field truncation effects are negligible for
$B/h_{\max}\ge40$.
The iteration counts remain nearly constant across the sweep, \chg{R1C3}{and the ratio
$t_{\mathrm{mesh}}/t_{\mathrm{solve}}$ increases moderately with problem size,
reflecting the expected growth of meshing cost relative to the nonlinear solve.}

\begin{table}[h!]
	\centering
	\caption{Far-field cap sensitivity at fixed interface resolution
		$n_{\mathrm{coh}}=80$ and $\theta_2=20^\circ$.}
	\label{tab:farfield_sweep}
	\begin{tabular}{rcccc}
		\toprule
		$B/h_{\max}$ & DOFs & $\sigma_{\mathrm{cr}}$ [MPa]
		& $t_{\mathrm{mesh}}/t_{\mathrm{solve}}$ & iters \\
		\midrule
		10 & 14{,}952  & 5.88913 & 0.72 & 18 \\
		20 & 21{,}216  & 5.87796 & 0.59 & 18 \\
		40 & 45{,}356  & 5.87548 & 0.63 & 17 \\
		80 & 133{,}184 & 5.87536 & 1.14 & 17 \\
		\bottomrule
	\end{tabular}
\end{table}

Based on these results, production simulations adopt
$n_{\mathrm{coh}}\ge40$ together with $B/h_{\max}\ge40$.
Under these conditions, the critical load is converged to within
$\epsilon\lesssim10^{-6}$, and further refinement of either the cohesive
discretization or the far-field truncation does not affect the results in a
meaningful way.


\subsection{Selection of the kink angle}
\label{sec:kink-angle}

\chg{R3C3}{
	\paragraph{On the definition of the critical state and the role of selection criteria.}
	In the present framework, the critical state is defined uniquely and
	independently of any directional selection.
	For each prescribed trial orientation of the cohesive continuation, the
	critical state is reached when the normal separation at the mouth of the
	cohesive zone attains its critical value, $d_n = d_n^{\max}$, corresponding to
	the exhaustion of normal cohesive resistance.
	The criteria introduced below do not redefine this failure condition.
	Instead, they provide alternative \emph{selection measures} for identifying
	the preferred kink direction, all evaluated at the same critical state.
	Accordingly, they should be interpreted as complementary diagnostics probing
	different aspects of the critical configuration rather than as competing
	failure criteria.
}

\paragraph{Load-based selector.}
For a trial direction $\theta_2$, the critical state is solved and the
corresponding critical load $\sigma_{\mathrm{cr}}(\theta_2)$ is recorded.
The kink angle is defined as
\[
\theta_2^{\sigma}
=
\operatorname*{arg\,min}_{\theta_2 \in \Theta}
\ \sigma_{\mathrm{cr}}(\theta_2),
\]
where $\Theta$ denotes the prescribed search set of global angles.
Numerically, the minimum is first localized on a coarse sweep and subsequently
refined on a fine sweep; the reported minimizer corresponds to the vertex of a
three--point parabolic interpolation through the discrete minimum and its two
nearest neighbors.

With $(n_{\mathrm{coh}}, B/h_{\max}) = (40,40)$ and $h_{\mathrm{grad}}=1.15$, the
coarse sweep yields
$\min \sigma_{\mathrm{cr}} = 5.65160\,\text{MPa}$ at
$\theta_2 \simeq 1.833^\circ$.
Refinement to $(n_{\mathrm{coh}}, B/h_{\max}) = (80,80)$ gives
$\min \sigma_{\mathrm{cr}} = 5.65151\,\text{MPa}$ at
$\theta_2 \simeq 1.727^\circ$,
corresponding to an angular shift of $0.106^\circ$.
The corresponding coarse and refined curves are shown in
Fig.~\ref{fig:kink-sweeps}, where diamonds mark the parabolic vertex and dashed
lines indicate its position.

\begin{figure}[t]
	\centering
	\begin{subfigure}[t]{0.48\linewidth}\centering
		\includegraphics[width=\linewidth]{Fig6aNew}
		\caption{Coarse sweep ($n_{\mathrm{coh}}=40$, $B/h_{\max}=40$).}
		\label{fig:kink-sweeps:a}
	\end{subfigure}\hfill
	\begin{subfigure}[t]{0.49\linewidth}\centering
		\includegraphics[width=\linewidth]{Fig6bNew}
		\caption{Fine sweep ($n_{\mathrm{coh}}=80$, $B/h_{\max}=80$).}
		\label{fig:kink-sweeps:b}
	\end{subfigure}
	\caption{Critical load $\sigma_{\mathrm{cr}}$ versus global cohesive-leg angle
		$\theta_2$. Diamonds denote the three--point parabolic vertex; dashed lines
		pass through the vertex.}
	\label{fig:kink-sweeps}
\end{figure}

\paragraph{Energetic selector.}
As a local alternative, we minimize the mixed-mode cohesive energy stored at
the mouth of Leg~2 in the critical state, normalized by the mode~I fracture
energy:
\begin{equation}
	\theta_2^{\Psi}
	=
	\operatorname*{arg\,min}_{\theta_2 \in \Theta}
	\bigl(\hat\Psi_{\mathrm{crit}}(\theta_2)-1\bigr),
	\qquad
	\hat\Psi_{\mathrm{crit}}=\Psi_{\mathrm{crit}}/\phi_n .
\end{equation}
Here \(\Psi_{\mathrm{crit}}\) denotes the mixed-mode cohesive potential evaluated
at the mouth at the onset of instability.

For the trapezoidal traction--separation laws employed in this work,
the normal opening at the mouth satisfies \(\bar d_n=1\) at the critical state,
so that \chg{R2C3}{\(F(\bar d_n=1)=c_n\)}. Substituting into the mixed-mode potential yields
\[
\hat\Psi_{\mathrm{crit}}
=
1
+
\left(
\frac{\phi_t}{\phi_n}\frac{1}{c_t}
-
r\sqrt{\frac{\phi_t}{\phi_n}}
\sqrt{\frac{c_n}{c_t}}
\right)
F(\bar d_t).
\]
The energetic excess \(\hat\Psi_{\mathrm{crit}}-1\) is therefore proportional to
the tangential cohesive potential \(F(\bar d_t)\). Minimizing this quantity is
thus equivalent to minimizing the shear opening at the mouth.


Applying the same two--stage sweep and parabolic vertex refinement strategy, the
coarse mesh yields
$\theta_2^{\Psi} \simeq 3.536^\circ$ with
$\min(\hat\Psi_{\mathrm{crit}}-1) \approx -1.876\times10^{-8}$.
The refined sweep gives
$\theta_2^{\Psi} \simeq 3.586^\circ$ with
$\min(\hat\Psi_{\mathrm{crit}}-1) \approx -6.131\times10^{-9}$,
corresponding to an angular shift of $0.05^\circ$.
The near--zero values confirm numerical consistency of the energetic diagnostic.
The corresponding results are shown in Fig.~\ref{fig:kink-psi}.

\begin{figure}[h!]
	\centering
	\begin{subfigure}[t]{0.48\textwidth}\centering
		\includegraphics[width=\linewidth]{Fig7aNew}
		\caption{Coarse sweep ($n_{\mathrm{coh}}=40$, $B/h_{\max}=40$).}
		\label{fig:kink-psi:a}
	\end{subfigure}\hfill
	\begin{subfigure}[t]{0.48\textwidth}\centering
		\includegraphics[width=\linewidth]{Fig7bNew}
		\caption{Fine sweep ($n_{\mathrm{coh}}=80$, $B/h_{\max}=80$).}
		\label{fig:kink-psi:b}
	\end{subfigure}
	\caption{Energetic mouth diagnostic
		$\hat\Psi_{\mathrm{crit}}-1$ versus global angle $\theta_2$.
		Diamonds indicate the parabolic vertex; dashed lines mark its position.}
	\label{fig:kink-psi}
\end{figure}


\paragraph{Critical-state fields.}
We illustrate the critical configuration for a representative coarse discretization
($n_{\mathrm{coh}}=40$, $h_{\mathrm{grad}}=1.15$, $B/h_{\max}=40$),
using the energetic selector $\theta_2 \approx 3.586^\circ$ for Leg~2.
Figure~\ref{fig:crit-results:stress} shows the normalized bulk stress component
$\sigma_{yy}/\sigma_n^{\max}$, exhibiting the expected near-interface concentration
and smooth spatial decay.
The underlying mesh is shown explicitly, confirming that the stress field remains
well resolved in the graded region adjacent to the cohesive leg.

Figure~\ref{fig:crit-results:profiles} reports profiles sampled along the cohesive leg
from the knee to the tip.
The top panel shows the normalized openings $\bar d_n$ and $\bar d_t$, indicating that
the normal separation is already close to its critical value near the knee, with
negligible tangential opening.
The middle and bottom panels compare the bulk tractions
(Cauchy stress projected onto $\pm\mathbf n$) on the upper and lower faces with the
TSL tractions reconstructed from the local displacement jump.

Excellent agreement is observed along most of the cohesive leg, \emph{excluding} a short
segment near the knee that lies on the softening branch of the cohesive law, with the
correct action--reaction symmetry (top $\bar\sigma_n \approx -1$, bottom
$\bar\sigma_n \approx +1$, and $\bar\sigma_t \approx 0$).
Small discrepancies are confined to the first few cohesive elements measured from the
knee and are associated with the early part of the softening branch.
These deviations are attributed to nodal stress recovery effects in the graded mesh
and do not affect the global consistency between bulk stresses and the imposed cohesive
traction--separation law.


\begin{figure}[!h]
	\centering
	\begin{subfigure}[t]{0.53\textwidth}
		\centering
		\includegraphics[width=\linewidth]{Fig8aNew}
		\caption{Normalized bulk $\sigma_{yy}$ near the cohesive leg (critical state). The deformed mesh is shown at 30× scale.}
		\label{fig:crit-results:stress}
	\end{subfigure}\hfill
	\begin{subfigure}[t]{0.45\textwidth}
		\centering
		\includegraphics[width=\linewidth]{Fig8bNew}
		\caption{Along-leg profiles: $\bar d_n,\bar d_t$ (top) and bulk vs.\ TSL tractions on the upper (middle) and lower (bottom) faces.}
		\label{fig:crit-results:profiles}
	\end{subfigure}
	\caption{Critical-state fields for the two-leg crack with a short cohesive continuation.
		Bulk tractions satisfy the cohesive boundary conditions and match the TSL prediction; minor
		differences are restricted to the first six cohesive elements from the knee.}
	\label{fig:crit-results}
\end{figure}

\chg{R1C4}{
	\paragraph{Conditioning of the load-based selector.}
	As illustrated in Fig.~\ref{fig:kink-sweeps}, the function
	$\sigma_{\mathrm{cr}}(\theta_2)$ exhibits a very shallow minimum in the
	vicinity of $\theta_2 \approx 1.6$--$1.9^\circ$.
	Although the minimum can be localized reliably using parabolic refinement in
	the present deterministic setting, the pronounced flatness of
	$\sigma_{\mathrm{cr}}(\theta_2)$ indicates that the load-based criterion is
	only weakly conditioned with respect to the kink angle.
	Consequently, in analyses affected by numerical noise (e.g., discretization
	errors or solver tolerances), or in materials characterized by a more
	extended fracture process zone, small perturbations may lead to noticeable
	variations in the estimated minimizing angle.
	This sensitivity motivates the use of complementary selection measures.
	In contrast to the load-based response, the energetic mouth diagnostic
	$\hat{\Psi}_{\mathrm{crit}}(\theta_2)-1$ exhibits a distinctly steeper and
	better-localized minimum, thereby providing improved numerical robustness
	when the global critical load is nearly insensitive to the kink direction.
}



\subsection{Discussion}
The two selectors probe different facets of the critical state and therefore
need not coincide.
The load-based criterion is \emph{global}, as it reflects equilibrium and the
mixed-mode response of the entire structure at the critical state.
In the present calculations, it yields
$\theta_2^{\sigma} \approx 1.6$--$1.9^\circ$, and the function
$\sigma_{\mathrm{cr}}(\theta_2)$ exhibits an extremely shallow minimum in this
vicinity.
As a result, the load-based selector is only weakly conditioned with respect
to the kink angle, and its minimizer is sensitive to discretization details
and solver tolerances.

By contrast, the energetic mouth criterion is inherently \emph{local}.
At the critical state, the normal opening at the mouth satisfies
$\bar d_n=1$, so that the normal cohesive response is already on its plateau
or early softening branch.
The energetic excess $\hat{\Psi}_{\mathrm{crit}}-1$ therefore scales with the
tangential contribution, which remains small and leads to a rapidly vanishing
value.
This produces a steeper and better-localized minimum of
$\hat{\Psi}_{\mathrm{crit}}(\theta_2)-1$, yielding
$\theta_2^{\Psi} \approx 3.586^\circ$.
The resulting difference between the two selectors is approximately
$1.86^\circ$.

In practice, the energetic diagnostic exhibits a pronounced U-shaped profile,
making the parabolic vertex fit robust even on relatively coarse angular
sweeps.
In contrast, the much flatter minimum of $\sigma_{\mathrm{cr}}(\theta_2)$
benefits significantly from refinement and tight bracketing of the search
interval.
If a single selector is desired, a natural compromise is to minimize the
\emph{integrated} cohesive work along Leg~2 rather than its mouth value alone,
or to employ a composite objective that balances load capacity against cohesive
dissipation, thereby accounting for the response of the full interface.


\chg{R3C1}{
	\paragraph{Remark on cohesive-law regularization and modeling scope.}
	The interpretation of the above results is closely tied to the adopted cohesive
	regularization.
	The present framework employs a smooth, potential-based traction--separation
	law, with the displacement jump as the primary kinematic variable and cohesive
	tractions obtained as its energetic conjugate.
	This choice is motivated by numerical robustness and the availability of a
	symmetric algorithmic tangent in mixed-mode problems.
	Within this setting, the extent of the \emph{active} cohesive zone is not
	prescribed a priori, but emerges implicitly from the solution as the portion of
	the interface undergoing separation.
	Formulations in which the cohesive-zone length is treated as an explicit
	unknown correspond to free-boundary cohesive crack models and require a
	different mathematical setting, which is beyond the scope of the present
	finite-element CZM approach.
}

\chg{R3C2}{
	\paragraph{Remark on the role of crack and trial-leg lengths.}
	Consistent with this modeling scope, both the initial crack and the kinked
	segment are introduced here as given geometrical inputs in the incipient-kinking
	analysis.
	The length of the kinked segment should be regarded as an auxiliary numerical
	parameter rather than as a physical crack length.
	The physically relevant quantity is the \emph{active} cohesive-zone extent,
	which emerges from the solution as the portion of the trial interface that
	experiences separation.
	Accordingly, the trial leg is chosen sufficiently long so that the active
	cohesive zone remains strictly contained within it; once this consistency
	condition is satisfied, the predicted kink direction and critical load are
	insensitive to the selected trial-leg length.
	The estimate $\ell_p$ serves as a classical scaling guideline to ensure that
	the fracture-process zone can fully develop without artificial truncation; it
	does not prescribe the cohesive-zone length and is not a constitutive
	assumption.
	Treating the cohesive crack length as a true free parameter requires a
	free-boundary or incremental propagation formulation, which lies beyond the
	scope of the present incipient-kinking study.
}

\chg{R1C3}{
	\paragraph{Remark on a possible extension to crack propagation.}
	While the present study is restricted to the determination of the kinking
	direction and associated critical load for a predefined trial leg, the proposed
	framework naturally lends itself to extension toward incremental crack
	propagation.
	In such a setting, the crack path would be updated step by step, and a new
	cohesive segment would be activated at each increment.
	The thin-channel construction and subsequent collapse would then be repeated
	for the updated crack geometry, implying regeneration of the mesh only in the
	vicinity of the advancing crack.
	Although this introduces additional computational cost, the channel-based
	meshing procedure is fully automated, local to the crack region, and does not
	require explicit mesh surgery or global topological modifications; the
	interface operators and nonlinear solution strategy remain unchanged.
	The parametric kinking studies reported here indicate that the cost of mesh
	generation is comparable to that of the nonlinear solve and scales smoothly
	with interface refinement, suggesting that multi-step crack-growth simulations
	can be carried out without a prohibitive increase in computational effort.
}


\section{Comparison with LEFM predictors}
\label{sec:LEFM-comparison}

This section compares the crack-kinking direction obtained from the cohesive-zone
framework with the classical LEFM prediction.
The goal is to quantify how mode mixity, encoded in the orientation $\theta_{1}$ of the
physical leg, influences the preferred deflection angle and to assess whether the cohesive
continuation modifies the geometric kink condition.

\subsection{Relative angle definition}
Throughout the paper, the kink angle is measured as the signed relative rotation
between the two legs, consistent with the notation introduced in Table~\ref{tab:notation},
\begin{equation}
	\Delta\theta = \operatorname{wrap}_{(-\pi,\pi]}\!\left( \theta_{2} - \theta_{1} \right).
\end{equation}
This wrapping operation ensures that $\Delta\theta$ is always reported in the interval $(-\pi,\pi]$, preserving the correct direction of kinking (e.g., clockwise vs. counterclockwise). In plots, angles are reported in degrees, with $\Delta\theta$ obtained from the
wrapped value converted accordingly.

In the comparison below, we use two decorated forms of the same quantity:
\[
\Delta\theta_{\mathrm{LEFM}}, \qquad
\Delta\theta_{\mathrm{CZM}}
= \operatorname{wrap}_{(-\pi,\pi]}\!\left( \theta_{2} - \theta_{1} \right),
\]
where $\Delta\theta_{\mathrm{LEFM}}$ is obtained from the
maximum tangential stress (MTS) criterion, and $\Delta\theta_{\mathrm{CZM}}$ from the
cohesive critical state.


\subsection{LEFM predictor}
\label{sec:LEFM-predictor}

For a given inclination $\theta_{1}$ of the physical crack, the classical prediction
of the kinking direction within LEFM is obtained from the mixed-mode
stress intensity factors $(K_{\rm I},K_{\rm II})$ associated with
Leg~1 and the maximum--tangential--stress (MTS) criterion.  For inclined edge
cracks under uniform loading, closed-form approximations for
$(K_{\rm I},K_{\rm II})$ and the corresponding kink angle exist and provide
convenient engineering estimates.  In this work, however, all mixed-mode
quantities are computed directly by mode-separating the $J$-integral.

\paragraph{Single-contour $J$-integral and mirror-point mode decomposition.}
A circular contour of radius $r_{\rm I}$ is placed around the crack tip, and the
near-tip fields are sampled at pairs of mirror-symmetric points
\[
P(\theta) = \bigl(r_{\rm I}\cos\theta, \; r_{\rm I}\sin\theta\bigr),\qquad
Q(\theta) = \bigl(r_{\rm I}\cos\theta, \; -r_{\rm I}\sin\theta\bigr),
\qquad \theta\in(0,\pi).
\]
Let $\boldsymbol{\varepsilon}_{P}$, $\boldsymbol{\sigma}_{P}$ and
$\partial\mathbf{u}_{P}/\partial x_{1}$ denote the strain, stress, and directional
displacement gradient at~$P(\theta)$ in the local coordinates $(x_{1},x_{2})$ aligned
with the crack.  Their counterparts at $Q(\theta)$ are constructed analogously.
Following the mirror-point formulation of Ishikawa, Kitagawa and Okamura
(see Kuna~\cite{Kuna2013}), symmetric and antisymmetric combinations yield the
mode I and mode II components:
\[
\begin{aligned}
	\boldsymbol{\varepsilon}^{\rm I}
	&= \tfrac12\bigl[\varepsilon_{P,11} + \varepsilon_{Q,11},\;
	\varepsilon_{P,22} + \varepsilon_{Q,22},\;
	\gamma_{P,12} - \gamma_{Q,12}\bigr]^{T},\\[3pt]
	\boldsymbol{\varepsilon}^{\rm II}
	&= \tfrac12\bigl[\varepsilon_{P,11} - \varepsilon_{Q,11},\;
	\varepsilon_{P,22} - \varepsilon_{Q,22},\;
	\gamma_{P,12} + \gamma_{Q,12}\bigr]^{T},
\end{aligned}
\]
\[
\begin{aligned}
	\boldsymbol{\sigma}^{\rm I}
	&= \tfrac12\bigl[\sigma_{P,11} + \sigma_{Q,11},\;
	\sigma_{P,22} + \sigma_{Q,22},\;
	\sigma_{P,12} - \sigma_{Q,12}\bigr]^{T},\\[3pt]
	\boldsymbol{\sigma}^{\rm II}
	&= \tfrac12\bigl[\sigma_{P,11} - \sigma_{Q,11},\;
	\sigma_{P,22} - \sigma_{Q,22},\;
	\sigma_{P,12} + \sigma_{Q,12}\bigr]^{T}.
\end{aligned}
\]
Here $\gamma_{12}=2\varepsilon_{12}$ denotes the engineering shear strain.


The corresponding directional displacement gradients are
$\partial\mathbf{u}^{\rm I}/\partial x_{1}$ and
$\partial\mathbf{u}^{\rm II}/\partial x_{1}$.
The elastic energy densities are
\[
U_{\rm I} = \tfrac12\,\boldsymbol{\sigma}^{\rm I}\!:\boldsymbol{\varepsilon}^{\rm I},
\qquad
U_{\rm II} = \tfrac12\,\boldsymbol{\sigma}^{\rm II}\!:\boldsymbol{\varepsilon}^{\rm II}.
\]

\paragraph{Mode-separated $J$-integrals.}
Let $\mathbf{n}=(\cos\theta,\sin\theta)$ denote the outward normal at $P(\theta)$.
With $\mathbf{t}^{\rm I}=\boldsymbol{\sigma}^{\rm I}\mathbf{n}$ and
$\mathbf{t}^{\rm II}=\boldsymbol{\sigma}^{\rm II}\mathbf{n}$, the mode-separated
$J$-integrands projected onto~$x_{1}$ are
\[
f_{\rm I}(\theta)
= U_{\rm I} n_{1}
- \mathbf{t}^{\rm I}\!\cdot\!
\frac{\partial\mathbf{u}^{\rm I}}{\partial x_{1}},
\qquad
f_{\rm II}(\theta)
= U_{\rm II} n_{1}
- \mathbf{t}^{\rm II}\!\cdot\!
\frac{\partial\mathbf{u}^{\rm II}}{\partial x_{1}}.
\]
Integration over the upper semicircle and doubling for symmetry gives
\[
J_{\rm I}  = 2r_{\rm I}\!\int_{0}^{\pi} f_{\rm I}(\theta)\,d\theta,
\qquad
J_{\rm II} = 2r_{\rm I}\!\int_{0}^{\pi} f_{\rm II}(\theta)\,d\theta.
\]

\paragraph{Stress intensity factors and MTS kink angle.}
Using the plane-strain modulus $E' = E/(1-\nu^{2})$, the mode-separated
stress intensity factors follow from the energetic identities
\[
K_{\rm I} = \sqrt{E'J_{\rm I}},\qquad
K_{\rm II} = \operatorname{sign}(J_{\rm II})\sqrt{E'|J_{\rm II}|}.
\]
The LEFM kink angle is then obtained from the MTS condition
\[
\Delta\theta_{\rm LEFM}
= \arg\max_{\theta}\,\sigma_{\theta\theta}(K_{\rm I},K_{\rm II},\theta),
\]
where $\sigma_{\theta\theta}$ is the LEFM near-tip circumferential stress.

\paragraph{Critical load: Quadratic energetic criterion.}
For critical-load predictions, the mode-separated energy-release rates at a
reference load $\sigma_{0}$ are
\[
G_{\rm I}(\sigma_{0}) = \frac{K_{\rm I}^{2}}{E'}, \qquad
G_{\rm II}(\sigma_{0}) = \frac{K_{\rm II}^{2}}{E'}.
\]
Since $G\propto \sigma^{2}$ for fixed geometry,
\[
G_{\rm I}(\sigma)
= G_{\rm I}(\sigma_{0})\bigl(\tfrac{\sigma}{\sigma_{0}}\bigr)^{2},\qquad
G_{\rm II}(\sigma)
= G_{\rm II}(\sigma_{0})\bigl(\tfrac{\sigma}{\sigma_{0}}\bigr)^{2}.
\]
The quadratic mixed-mode criterion
\[
\frac{G_{\rm I}}{G_{\rm Ic}}
+ \frac{G_{\rm II}}{G_{\rm IIc}} = 1
\]
then yields the LEFM critical load
\[
\sigma_{\rm cr}^{\rm (quad)}
= \frac{\sigma_{0}}
{\sqrt{\dfrac{G_{\rm I}(\sigma_{0})}{G_{\rm Ic}}
		+ \dfrac{G_{\rm II}(\sigma_{0})}{G_{\rm IIc}}}}.
\]

\paragraph{Critical load: Benzeggagh--Kenane (BK) criterion.}
The BK criterion \cite{BenzeggaghKenane1996} is an energy-based mixed-mode
fracture model that prescribes an effective fracture toughness
\[
G_{\rm c}(\beta)
= G_{\rm Ic}
+ \bigl(G_{\rm IIc} - G_{\rm Ic}\bigr)\,\beta^{\eta},
\qquad
\beta = \frac{G_{\rm II}}{G_{\rm I}+G_{\rm II}},
\]
where $\eta>0$ controls the sensitivity of the toughness to the mode mixity.
Evaluating the mode-mix parameter $\beta_{0}$ at a reference load $\sigma_{0}$
allows the LEFM critical load to be computed from
$G(\sigma_{\rm cr}) = G_{\rm c}(\beta_{0})$, which gives
\[
\sigma_{\rm cr}^{\rm (BK)}(\eta)
= \sigma_{0}\,
\sqrt{\frac{G_{\rm c}(\beta_{0})}{
		G_{\rm I}(\sigma_{0}) + G_{\rm II}(\sigma_{0})}}.
\]
This expression provides a convenient mixed-mode benchmark for comparison with
the cohesive predictions.


\subsection{LEFM vs CZM mixed-mode fracture predictions}

In the cohesive-zone framework, each loading inclination $\theta_1$ is analyzed
through a two-stage procedure.
For each prescribed trial direction $\theta_2$ of the cohesive continuation,
a fully nonlinear equilibrium problem is solved up to the critical state.
An outer parametric sweep over $\theta_2$ is then performed, and the preferred
kink direction is selected according to a chosen diagnostic, either the global
critical stress $\sigma_{\mathrm{cr}}$ or the energetic mouth functional
$\hat{\Psi}$.
This approach yields consistent predictions for both the kink direction and the
associated critical load within a unified CZM setting.

Figure~\ref{fig:LEFM-vs-CZM} compares the resulting kink angles and critical
loads with classical LEFM predictions.
LEFM results are shown for the MTS criterion, as well as for quadratic and
Benzeggagh--Kenane (BK) mixed-mode laws with different $\eta$ parameters.
CZM results are reported for both selection strategies, namely minimization of
$\sigma_{\mathrm{cr}}$ and minimization of $\hat{\Psi}$.


\begin{figure}[h!]
	\centering
	\begin{subfigure}[b]{0.48\linewidth}
		\includegraphics[width=\linewidth]{Fig9aNew}
		\caption{Kink angle $\Delta\theta$ vs. $\theta_1$}
		\label{fig:kink-angle-subplot}
	\end{subfigure}
	\hfill
	\begin{subfigure}[b]{0.48\linewidth}
		\includegraphics[width=\linewidth]{Fig9bNew}
		\caption{Critical load $\sigma_{\rm cr}$ vs. $\theta_1$}
		\label{fig:crit-load-subplot}
	\end{subfigure}
	\caption{Comparison of LEFM and CZM predictions. LEFM includes MTS, quadratic, and BK criteria. CZM results correspond to minimization of either $\sigma_{\rm cr}$ or the energetic functional $\hat{\Psi}$.}
	\label{fig:LEFM-vs-CZM}
\end{figure}


\subsection*{Analysis of results}

\chg{R2C9}{
	\paragraph{Kink angle.}
	Figure~\ref{fig:kink-angle-subplot} compares the crack deflection angles predicted by LEFM and by CZM using two different selectors.	For small loading inclinations $\theta_1$, all approaches yield very close kink angles, with differences below $0.5^\circ$, indicating that the cohesive-zone formulation correctly recovers the LEFM limit when the fracture process zone is	small and the response is predominantly mode~I. \newline
	As $\theta_1$ increases, systematic differences emerge. The CZM prediction obtained by minimizing the critical load $\sigma_{\mathrm{cr}}$	consistently produces slightly \emph{shallower} deflection angles than LEFM, with deviations growing monotonically and reaching about $1.5^\circ$ at $\theta_1=30^\circ$. In contrast, the energetic selector based on minimizing $\hat{\Psi}$ yields systematically larger deflection magnitudes $|\Delta\theta|$ than the
	load-based selector, with deviations approaching $3^\circ$ at the largest
	inclination considered.\newline
	The smooth, nearly linear evolution of all three curves confirms that the
	outer minimization procedure is well posed.
	The larger separation observed for the energetic selector reflects its local
	character: by penalizing shear participation at the cohesive-zone mouth, it
	favors crack paths that promote opening-dominated separation even when the
	global load capacity varies weakly with direction.
}


\paragraph{Critical load.}
The corresponding critical loads are shown in
Fig.~\ref{fig:crit-load-subplot}.
For all loading angles $\theta_1$, the CZM predicts lower critical stresses than
LEFM, with the difference increasing as mode~II participation grows.
At $\theta_1=2^\circ$, the reduction is below $1.5\%$, whereas at
$\theta_1=30^\circ$ it reaches approximately $10\%$.

The two CZM selectors produce virtually identical values of
$\sigma_{\mathrm{cr}}$, with differences remaining below $0.05\%$ across the
entire range of $\theta_1$.
This confirms that, although the selectors may lead to different kink
directions, the global load-carrying capacity is governed primarily by the
overall cohesive dissipation rather than by the local directional criterion.


\subsection*{Physical Interpretation}

The parameter $\theta_1$ controls the mode mixity at the crack tip rather than the
applied load level.
For small $\theta_1$, the response is opening-dominated and LEFM and CZM yield
nearly identical kink angles and critical loads, reflecting convergence of the
cohesive formulation to the LEFM limit.
As $\theta_1$ increases, shear participation modifies the distribution of
tractions and separations within the fracture process zone, even though its
overall extent remains governed primarily by material parameters and the
attained load.
This internal redistribution is captured by the mixed-mode cohesive law without
invoking any notable change in process-zone length.
The distinction between the CZM selectors reflects their different physical
emphases: minimization of the global critical load $\sigma_{\mathrm{cr}}$
identifies directions close to the LEFM prediction, whereas minimization of the
energetic mouth measure $\hat{\Psi}$ emphasizes local opening at the knee of the
cohesive continuation, leading to systematically steeper kink angles while
leaving the critical load essentially unchanged.

\section{Sensitivity to the auxiliary cohesive-leg length, range of applicability, and limitations}
\label{sec:scope_limitations}

The critical state in the present formulation remains defined exactly as in
Section~\ref{sec:results}: for each prescribed trial orientation of
Leg~2, the nonlinear problem is solved until the mouth opening reaches the
critical normal separation. What is examined here is a different issue,
namely the influence of the \emph{auxiliary} cohesive-leg length~$\delta$ on
the resulting directional prediction and on the interpretation of the method.
This distinction is important because the present crack-kinking procedure does
not operate at a mathematical crack tip. Instead, it infers the preferred
direction from a finite cohesive continuation that regularizes the singular
elastic field over a short but nonzero process-zone probe.

The present sensitivity study is carried out in a deliberately near-LEFM
regime, taking $\beta/a \approx 0.05$ as a practical benchmark for a small
but finite fracture-process zone. This value should be understood as a
working reference for controlled comparison, not as a universal criterion for
the validity of LEFM.

\subsection{Calibrated adaptive estimation of the auxiliary cohesive length and its numerical robustness}
\label{subsec:adaptive_delta}

The present crack-kinking procedure does not evaluate a tip criterion at a
mathematical crack tip. Instead, it infers the preferred propagation direction
from a short cohesive continuation of finite length~$\delta$, which
regularizes the singular elastic field over a finite process-zone probe. For
this reason, the choice of~$\delta$ is not a secondary implementation detail:
it must be large enough to contain the full active cohesive zone, but it
should not vary in an uncontrolled way when material parameters are changed.

In the present formulation, $\delta$ is not treated as a material parameter.
It is an auxiliary numerical length introduced to represent the finite
fracture-process zone in a controlled way. If $\delta$ were kept fixed while
the cohesive strength varied, the activated part of the cohesive segment could
occupy very different fractions of the trial leg. In that case, changes in the
predicted kink direction would reflect not only the constitutive variation
itself, but also an inconsistent relation between the finite cohesive zone and
the imposed probe geometry. To avoid this difficulty, the present calculations
use a calibrated adaptive estimate of~$\delta$.

We introduce the stiffness-to-strength ratio
\begin{equation}
	\rho=\frac{E}{\sigma_n^{\max}},
\end{equation}
and define the effective stress-intensity measure
\begin{equation}
	K_{\mathrm{eff}}=\sqrt{K_I^2+K_{II}^2}.
\end{equation}
A Dugdale-type estimate is then used for the characteristic process-zone size,
\begin{equation}
	r_D=C_D\left(\frac{K_{\mathrm{eff}}}{\sigma_n^{\max}}\right)^2,
\end{equation}
where $C_D$ is a dimensionless calibration constant. The trial cohesive length
is chosen from
\begin{equation}
	\delta=\frac{r_D}{r_\beta},
\end{equation}
where $r_\beta$ is the target fraction of the active cohesive length within
the trial segment, and $\beta$ denotes the activated part of Leg~2 in the
converged critical state.

Here $C_D$ is not taken as universal. Rather, it is calibrated separately for
each $\rho$-sweep so that the activated cohesive zone occupies approximately
the same fraction of the trial segment throughout that sweep. For the three
initial crack inclinations considered here, the adopted values are
$C_D=17.2$ for $\theta_1=-30^\circ$,
$C_D=22.36=1.3\times 17.2$ for $\theta_1=-40^\circ$, and
$C_D=34.4=2.0\times 17.2$ for $\theta_1=-50^\circ$. Accordingly, the present
estimate of~$\delta$ should be interpreted as a calibrated predictor rather
than as a universal closed-form law.

This gives the procedure a predictor--corrector character. The LEFM field
provides a first estimate of the process-zone scale through $K_{\mathrm{eff}}$,
while the subsequent CZM solution checks a posteriori whether the activated
length~$\beta$ indeed occupies the intended fraction of the trial cohesive
segment. The role of the study below is therefore not merely parametric. It is
to verify that this normalization remains effective as the stiffness-to-strength
ratio changes.

A further algorithmic distinction should be emphasized. In Section~7, the kink
angle is presented through coarse and fine angular sweeps in order to
illustrate directly the shallow minimum of $\sigma_{cr}(\theta_2)$ and the
sharper minimum of the energetic selector. In the present parametric study,
however, the selected angle $\theta_2^\ast$ is determined by an automated
iterative procedure that repeatedly brackets and refines the minimizing angle
until the prescribed angular tolerance is reached. This replacement is purely
algorithmic: it does not alter the directional criterion itself, but makes the
determination of $\theta_2^\ast$ fully automatic and computationally efficient
for repeated calculations. In this sense, the procedure is already suitable as
the directional core of a future stepwise crack-growth algorithm, although the
corresponding propagation study lies beyond the scope of the present paper.

Table~\ref{tab:rho_sweep_all} summarizes the results obtained on the finer mesh
used for production calculations, with
\[
n_{\mathrm{coh}}=80,\qquad B/h_{\max}=80,\qquad h_{\mathrm{grad}}=1.15.
\]
Three conclusions follow immediately. First, the normalized trial length
$\delta/a$ increases monotonically with $\rho$ for all three initial crack
orientations. This is consistent with the estimate
$r_D\propto K_{\mathrm{eff}}^2/(\sigma_n^{\max})^2$: with fixed $E$,
increasing $\rho$ corresponds to decreasing cohesive strength and therefore to
a larger predicted process-zone size. Second, the ratio $\beta/\delta$ remains
nearly constant throughout each sweep. This confirms that the calibrated
LEFM-based estimate of~$\delta$ regularizes the finite cohesive problem in a
quantitatively controlled manner. Third, the selected kink angle
$\theta_2^\ast$ varies only weakly with~$\rho$ within each sweep. Thus, once
the process zone is represented consistently, moderate changes in the
stiffness-to-strength ratio primarily rescale the finite cohesive length while
affecting the directional prediction only slightly.

\begin{table}[t]
	\centering
	\small
	\setlength{\tabcolsep}{4pt}
	\caption{Results of the $\rho$-sweep with calibrated adaptive estimation of $\delta$
		for different initial crack orientations~$\theta_1$. Here $\theta_2^\ast$
		denotes the selected CZM kink angle.}
	\label{tab:rho_sweep_all}
	\begin{tabular}{c ccc ccc ccc}
		\toprule
		& \multicolumn{3}{c}{$\theta_1=-30^\circ$}
		& \multicolumn{3}{c}{$\theta_1=-40^\circ$}
		& \multicolumn{3}{c}{$\theta_1=-50^\circ$} \\
		\cmidrule(lr){2-4}\cmidrule(lr){5-7}\cmidrule(lr){8-10}
		$\rho$
		& $\delta/a$ & $\beta/\delta$ & $\theta_2^\ast$
		& $\delta/a$ & $\beta/\delta$ & $\theta_2^\ast$
		& $\delta/a$ & $\beta/\delta$ & $\theta_2^\ast$ \\
		\midrule
		80  & 0.0307 & 0.9000 & 4.4568 & 0.0300 & 0.9125 & 4.5000 & 0.0303 & 0.9125 & 3.3304 \\
		85  & 0.0347 & 0.9000 & 4.4568 & 0.0333 & 0.9375 & 4.5230 & 0.0342 & 0.9250 & 3.3475 \\
		90  & 0.0389 & 0.9000 & 4.4435 & 0.0373 & 0.9375 & 4.5154 & 0.0383 & 0.9250 & 3.3887 \\
		95  & 0.0433 & 0.9000 & 4.4435 & 0.0416 & 0.9375 & 4.5230 & 0.0427 & 0.9250 & 3.4037 \\
		100 & 0.0480 & 0.9000 & 4.4298 & 0.0461 & 0.9375 & 4.5358 & 0.0473 & 0.9250 & 3.4487 \\
		105 & 0.0529 & 0.9000 & 4.4298 & 0.0508 & 0.9375 & 4.5578 & 0.0522 & 0.9250 & 3.5248 \\
		110 & 0.0580 & 0.9000 & 4.4135 & 0.0558 & 0.9375 & 4.5200 & 0.0572 & 0.9250 & 3.5267 \\
		115 & 0.0634 & 0.9000 & 4.4135 & 0.0610 & 0.9375 & 4.5476 & 0.0626 & 0.9375 & 3.5773 \\
		120 & 0.0691 & 0.9000 & 4.4034 & 0.0664 & 0.9375 & 4.5831 & 0.0681 & 0.9375 & 3.6285 \\
		\bottomrule
	\end{tabular}
\end{table}


The robustness of these trends was checked by repeating the same
\(\rho\)-sweeps on the coarser mesh
\[
n_{\mathrm{coh}}=40,\qquad B/h_{\max}=40,\qquad h_{\mathrm{grad}}=1.15,
\]
and comparing the results with the production calculations reported in
Table~\ref{tab:rho_sweep_all}, obtained on the finer mesh
\[
n_{\mathrm{coh}}=80,\qquad B/h_{\max}=80,\qquad h_{\mathrm{grad}}=1.15.
\]
This comparison shows that \(\delta/a\) remains unchanged to the reported
precision for all parameter points. The ratio \(\beta/\delta\) is likewise
unchanged except for a single isolated difference of \(-0.0125\) at
\(\theta_1=-50^\circ\), \(\rho=115\). The corresponding changes in the
predicted angle \(\theta_2^\ast\) are small throughout: the maximum absolute
difference is \(0.0064^\circ\) for \(\theta_1=-30^\circ\),
\(0.0069^\circ\) for \(\theta_1=-40^\circ\), and \(0.0311^\circ\) for
\(\theta_1=-50^\circ\). These differences are much smaller than the total
variation of \(\theta_2^\ast\) across each \(\rho\)-sweep and do not modify
any of the trends discussed above. The weak dependence of the selected kink
angle on \(\rho\) in the present range therefore cannot be attributed to the
particular mesh used for the production runs.

These observations do not imply that the kinking direction is generally
independent of cohesive parameters. They show only that, within the present
quasi-static mixed-mode setting and for the class of cohesive laws considered
here, the calibrated predictor--corrector choice of the auxiliary cohesive
length provides a numerically stable basis for comparing directional
predictions across a family of parameter values.




\subsection{Relation to LEFM, applicability, and potential failure modes}
\label{subsec:scope_failure_modes}

The present results should be interpreted relative to LEFM with care. In
classical LEFM, the propagation criterion is evaluated at a mathematical crack
tip and depends only on asymptotic quantities such as $K_I$ and $K_{II}$. In
the present CZM framework, by contrast, the propagation direction is inferred
from a finite cohesive segment carrying distributed tractions. For this
reason, exact agreement with LEFM criteria is not expected in general, even
when the trends are similar. The two descriptions coincide only in an
asymptotic regime in which the process zone becomes vanishingly small
compared with all geometric scales. That limit is not established here.
Accordingly, the present method should be viewed as a finite-process-zone
formulation whose relation to LEFM is asymptotic rather than exact.

Within that interpretation, the method is expected to be most reliable for
quasi-static mixed-mode crack growth in quasi-brittle solids when the
fracture-process zone remains small compared with the ligament size and the
other structural dimensions. Outside this regime, several failure modes should
be anticipated. If the traction--separation law departs strongly from the
present class of laws, for example through pronounced tension--shear
asymmetry, strong rate dependence, or markedly different softening structure,
the Dugdale-type predictor for $\delta$ may cease to keep $\beta/\delta$
approximately constant and may require a different calibration or an
additional adaptive loop. If the process zone becomes comparable with the
structural dimensions, the local interpretation of the cohesive segment as a
short directional probe loses validity, and the predicted direction may then
reflect the imposed probe geometry as much as the underlying fracture
physics. If the interface discretization is insufficient, both the activated
length~$\beta$ and the directional diagnostic itself may be distorted. The
adaptive estimate of $\delta$ mitigates this difficulty but does not remove
the need for mesh verification. Finally, the present formulation does not
include inertia, large-scale plasticity, or fully developed viscoelastic
fracture; dynamic fracture, highly ductile response, and hereditary
constitutive effects therefore lie outside the validated scope of the present
study.

Taken together, these observations define the meaning of the present
validation more precisely. The results do not establish universal agreement
with LEFM, nor do they show parameter-independence of the CZM direction in a
general sense. They show, rather, that within the targeted quasi-static
quasi-brittle regime, the predictor--corrector choice of the auxiliary
cohesive length provides a consistent and numerically robust basis for
comparing directional predictions across a family of cohesive parameters.

\section{Conclusions}
\label{sec:conclusions}

This work developed and analysed a fully implicit cohesive-zone framework for
mixed-mode crack kinking.
 The formulation augments a pre-existing
crack by a short, zero-thickness cohesive extension, which enables a systematic
and well-posed evaluation of the critical state associated with any prescribed
trial kink orientation. Central to the approach is a potential-based mixed-mode
traction--separation law (TSL), which guarantees energetic consistency, a
symmetric algorithmic tangent, and robust quadratic convergence of the Newton
iterations. By combining this constitutive structure with a channel-collapse
meshing strategy, the framework provides a controlled numerical setting for
assessing mixed-mode kinking without invoking enrichment techniques,
geometry-specific assumptions, or explicit mesh surgery, and allows a direct
comparison between cohesive and elastic predictions.

\chg{R3C4}{
	The cohesive-zone model employed in this study is not introduced as a competing
	fracture theory nor as a material-specific law. Instead, it is used as a
	regularized and energetically consistent framework for assessing the robustness
	of classical LEFM-based kinking predictions in the presence of a finite but
	small fracture process zone. Within this scope, LEFM is interpreted as the
	asymptotic limit corresponding to vanishing cohesive-zone length, and the
	present analysis is conducted deliberately in a near-LEFM regime that permits
	controlled and meaningful comparison.
}

A comparison with classical LEFM-based kinking criteria leads to four main
observations.

\begin{enumerate}
	\item \emph{Kink direction is essentially an elastic quantity.}
	\chg{R1C5}{
		Across all loading inclinations investigated, the predicted cohesive kink
		angles remain within approximately one degree of the LEFM
		maximum-tangential-stress (MTS) prediction.
		This confirms the well-established LEFM result that the crack-deflection
		direction is governed primarily by the elastic mixed-mode stress field.
		The present study does not claim this as a new physical finding; rather, it
		demonstrates that the proposed cohesive-zone framework is capable of
		reproducing this classical behaviour in a numerically robust and energetically consistent
		manner.
	}
	In this sense, the cohesive process zone and the specific shape of the TSL have
	a negligible influence on the deflection angle, while their role is expressed
	predominantly through the associated critical load.
	
	\item \emph{Critical loads carry a clear imprint of process-zone physics.}
	Although both LEFM and the cohesive framework predict a monotonic increase of
	the critical load with crack inclination, the cohesive model consistently
	yields lower values. The reduction is modest (below $2\%$) near mode~I, but
	reaches approximately $10\%$ for strongly skewed cracks. This systematic
	difference reflects the finite size of the fracture-process zone and the
	nonlinear separation response encoded in the TSL—effects that are absent in
	classical LEFM.
	
	\item \emph{The potential-based cohesive law exhibits stronger mode interaction
		than BK-type criteria.}
	Comparisons with Benzeggagh--Kenane mixed-mode surfaces ($\eta=1$--$3$) show
	that all BK predictions overestimate the effective toughness relative to the
	present cohesive model. The cohesive critical loads lie below all BK curves,
	indicating a stronger energetic coupling between modes, which effectively
	corresponds to a mixed-mode exponent $\eta_{\mathrm{eff}}>3$. This highlights
	that the potential-based TSL introduces mode interaction beyond that captured
	by standard empirical criteria.
	
	\item \emph{The numerical framework shows high robustness and reproducibility.}
	Convergence studies confirm that accurate and mesh-independent critical loads
	are obtained once two modest conditions are met: (i) a bulk refinement level
	$B/h_{\max}\ge40$, and (ii) at least $n_{\mathrm{coh}}\ge40$ cohesive segments
	(or twice this value for conservative runs). Under these conditions, the
	relative error in the critical load falls below $10^{-6}$, and reconstructed
	bulk tractions closely match the cohesive tractions along nearly the entire
	cohesive leg. Among the two tested selection measures, the global minimum of
	$\sigma_{\mathrm{cr}}(\theta_2)$ is shallow and sensitive to bracketing, whereas
	the energetic mouth criterion exhibits a much sharper minimum and is
	numerically more robust.
\end{enumerate}

\chg{R1C5}{
	Overall, the principal contribution of this study lies not in proposing a new
	fracture criterion or physical crack-deflection mechanism, but in establishing
	a cohesive-zone formulation that is energetically consistent, numerically
	robust, and capable of assessing the robustness of classical LEFM kinking
	predictions within a regularized, non-singular framework. The unified cohesive
	potential, symmetric Hessian tangent, and channel-collapse meshing strategy
	together provide a stable and reproducible tool for determining both kink
	direction and critical load in mixed-mode cohesive-zone analyses.
}

\paragraph{Outlook: Sensitivity to TSL shape.}
Because the predicted critical load depends on details of the
traction--separation law—including strengthening length, pla\-teau width, and
softening curvature—a broader parametric investigation of TSL shapes is
warranted. Earlier mode-I studies (e.g.,~\cite{Selivanov2019ModeI}) indicate that
the structure of the fracture process zone can be highly sensitive to these
features even when global response quantities vary only moderately. Extending
such analyses to fully mixed-mode kinking represents a natural next step. The
present framework can also be extended to incremental crack growth by
reorienting the cohesive leg as the crack advances, thereby providing a
consistent starting point for future propagation studies.




\printbibliography



\end{document}

